% Source-faithful fragments for §§7–11. Generated by build_geometry.py.
\section{Automorphisms and real structures}
\label{sec:real}
The full holomorphic automorphism group is the theorem of
\cite[theorem \texttt{p2:thm:full-automorphisms}]{DGKXv1}; its identity
component and translation action are due to \cite[\S9.7, Proposition~9.23
and Corollary~9.24]{AC26}. Our lattice calculation is a re-derivation.
\begin{cited}[FAM B.4.1; K]\label{fam:B.4.1}
For each member of the retained $(3,4)$ family, $\operatorname{Aut}(X_c)$, including all its components, is $\mathbb C/\mathbb Z\simeq\mathbb C^*$. Its action on the ordinary fibres is $\zeta\mapsto\zeta+be$, and the exponential group isomorphism is $b\mapsto e^{2\pi ib}$.
\end{cited}
\noindent\emph{Uses:} FAM A.1.1, A.1.2, A.2.1, A.3.1 and the cusp valuation comparison in FAM B.2.\par
\noindent The re-derivation is in Appendix~B.\par

\begin{proposition}[FAM B.4.2; T]\label{fam:B.4.2}
Use the normalization $(\beta^{\rm part}+\tau)(p_0)=0$. There is an antiholomorphic diffeomorphism $X_c\to X_{c'}$ exactly when
\[
 c'-(1-\bar c)\in2\mathbb Z,
 \qquad\text{equivalently }w'=-\bar w.
\]
Consequently $\overline{X_c}\simeq X_{1-\bar c}$. The members admitting antiholomorphic self-maps are exactly those with $\operatorname{Re}c\in\tfrac12+\mathbb Z$, equivalently $w=-\bar w$. On each such member the group of holomorphic and antiholomorphic diffeomorphisms is
\[
 \operatorname{Aut}^{\pm}(X_c)\simeq\mathbb C^*\rtimes\mathbb Z/2,
 \qquad z\longmapsto1/\bar z
\]
for the action of the second factor. Its antiholomorphic maps form the nonidentity coset. For all other members the full group is the holomorphic group of \ref{fam:B.4.1}.
\end{proposition}
\noindent\emph{Uses:} FAM A.3.1, A.5.1, B.2.1 and B.4.1; the cusp-normalized period laws.\par

\begin{proof}
\par\noindent\textit{Step 1.} We first construct the comparison before returning its affine data to the canonical choice. Lift conjugation of the real orbifold base to the antiholomorphic involution $\iota$ of its universal cover whose cusp branch has $\tau(\iota z)=-\overline{\tau(z)}$. A lift exists by the universal covering property; it extends at the finite orbifold points, whose local equations are $s_j^{m_j}=t_j$. Choose the usual paths from a real point $t>1$, so reflection fixes the second centre and sends the first chosen lift to $g_2^{-1}z_1$. The induced involution of the deck group is
\[
 \phi(g_1)=g_2^{-1}g_1^{-1}g_2,\qquad
 \phi(g_2)=g_2^{-1},\qquad \phi(g_0)=g_0^{-1}.
\]
These identities also follow by reflecting the lifted peripheral paths; their product and their squares agree with $g_1g_2g_0=1$. The cusp branch fixes the lift uniquely.

\par\noindent\textit{Step 2.} All reflection identities for the periods can be checked before invoking uniqueness. Denote their affine transformation formulas by $G_1,G_2$ on the triple $(\tau,\mu,\beta)$ and put $K(\tau,\mu,\beta)=(-\bar\tau,1-\bar\mu,-\bar\beta)$. Direct rational substitution gives
\[
 KG_1K=G_2^{-1}G_1^{-1}G_2,\qquad KG_2K=G_2^{-1}.
\]
For example $G_2^{-1}(\tau,\mu,\beta)=(-1/\tau,(1-\mu)/\tau,\beta+3-6(1-\mu)^2/\tau)$ and
$G_1^{-1}(\tau,\mu,\beta)=(1/(1-\tau),1-\mu/(1-\tau),\beta-2+6\mu^2/(1-\tau))$. These formulas verify every constant in the two equalities. The reflected $\mu$ is cusp bounded and has the same affine laws, so their difference is a global section of $\mathcal O_B(-p_0)$ and vanishes. After this equality, the difference between the reflected and original $\beta^{\rm part}$ is a holomorphic invariant base function, hence constant; its cusp value is zero under the stated normalization. Therefore
\[
 \mu(\iota z)=1-\bar\mu(z),\qquad
 \beta^{\rm part}(\iota z)=-\overline{\beta^{\rm part}(z)}.
\]
The original $\tau$ identity follows from its $j$ equation and its specified cusp branch. In particular $D_{-\bar c}(\iota z)=D_c(z)$, so this comparison stays in the exact admissible domain.

\par\noindent\textit{Step 3.} Let
\[
 H=\begin{pmatrix}1&0&0&0\\0&1&0&0\\-6&0&-1&0\\0&-1&0&-1\end{pmatrix}.
\]
Then $H^2=I$, $HA_2H=A_2^{-1}$, $HA_1H=A_2^{-1}A_1^{-1}A_2$, and
\[
 \Pi_{-\bar c}(\iota z)H=-\overline{\Pi_c(z)}.
\]
Thus $(z,\zeta)\mapsto(\iota z,-\bar\zeta)$ is a well-defined antiholomorphic torus map. It also descends under the middle deck action. Indeed covariance and the displayed lattice conjugacies imply $R_{\phi(g)}(\iota z)=\overline{R_g(z)}$ by multiplying the full real-invertible period matrices. This proves exact compatibility, not merely an equality of fibre lattices.

\par\noindent\textit{Step 4.} At the first finite centre compose this map with the target deck identification $g_2$ to return the centre $g_2^{-1}z_1$ to $z_1$. Its lattice map becomes $A_2H$ and it conjugates the clockwise generator to the inverse generator. At the second centre its lattice map is $H$ with the same reversal. Consequently the target affine vectors are exactly
\[
 w'_1=-A_2H\varepsilon=-\varepsilon+d,\qquad
 w'_2=-H(-\varepsilon')=\varepsilon'-3d.
\]
The minus signs arise because the inverse of a generator with invariant translation $w'_j/m_j$ has translation $-w'_j/m_j$. This calculation verifies the full affine actions on the root covers. For the overlaps choose the root coordinates of the real local base coordinates $t$ and $t-1$ compatible with these local reflections. This choice has an explicit comparison with any original root coordinate: their ratio is a holomorphic unit, so changing it changes the collar by $(2\pi i)^{-1}\log h\,\Pi w_j$, a translation extending on the entire root family. Its inverse gives a holomorphic isomorphism of the two glued descriptions. Thus the extension property and final classification transfer to the original coordinates. Then $\log s'_j=\overline{\log s_j}$ on paired branches. Under $\zeta\mapsto-\bar\zeta$, conjugation of $1/(2\pi i)$ supplies the compensating sign, and the identity of the two logarithmic collar translations follows exactly. Two roots of the same fixed real coordinate differ by an $m_j$th root of unity; this changes the comparison only by a constant period-coordinate translation, and the additive correction of \ref{fam:A.3.1} glues those constants. The preceding holomorphic-unit comparison handles a different original local base coordinate. Its inverse has the same construction.

\par\noindent\textit{Step 5.} At the cusp choose $s=(2\pi i)^{-1}\log t_c$, so $s\mapsto-\bar s$. The full holomorphic unit matrix satisfies
\[
 C_{-\bar c}(\bar t_c)=-\overline{C_c(t_c)}+
       \begin{pmatrix}6&0\\0&1\end{pmatrix}.
\]
The integral extra matrix has trivial exponential. Hence $(x,t_c)\mapsto(\bar x,\bar t_c)$ intertwines the full cusp deck groups, with the same integral index $\lambda$. The toric fan is defined over $\mathbb Z$, so this map extends antiholomorphically on every toric chart and stratum. It preserves the untwisted cusp collar. We have therefore constructed an antiholomorphic map to the actual target datum with parameter $-\bar c$, $r'_1=-1,k'_1=1,r'_2=1,k'_2=-3,l'_0=0$, whose $p'$ is $+1$.

\par\noindent\textit{Step 6.} Apply \ref{fam:B.2.1} to normalize this target back to the canonical $p=-1$ datum. Take $a=-1$ and the odd shift $n=1$. This gives a holomorphic map to $X_{1-\bar c}$. Every other antiholomorphic map, composed with the inverse of this one, is a holomorphic comparison between canonical members. \ref{fam:A.5.1} gives exactly the stated congruence. Exponentiating it gives $w'=-\bar w$.

\par\noindent\textit{Step 7.} It remains to verify the group extension rather than infer its square from $H^2$. For a real member, choose any resulting antiholomorphic self-map $J$. Its derivative on ordinary fibre coordinates is $+\text{complex conjugation}$: the initial reflection has derivative $-\text{complex conjugation}$, while the normalization just used has linear part $-I_2$. Hence it conjugates translation by $b\in\mathbb C/\mathbb Z$ to translation by $\bar b$. Its square is a holomorphic translation $d$. Commuting $J$ with $J^2$ shows $d=\bar d$ in $\mathbb C/\mathbb Z$. This equality forces the imaginary part of a representative of $d$ to vanish, so take $d\in\mathbb R$. Replacing $J$ by the composition of translation by $-d/2$ with $J$ changes its square to $d-d/2-d/2=0$. We have an antiholomorphic involution. Under $b\mapsto e^{2\pi ib}$, its conjugation action is $z\mapsto1/\bar z$. This proves the split group and specifies its antiholomorphic coset.
\end{proof}

\begin{theorem}[ARI C.2.2; T]\label{ari:C.2.2}
An antiholomorphic equivalence $X_c\to X_{c'}$ exists exactly when
$c'-(1-\bar c)\in2\mathbb Z$. Consequently $X_c$ admits an
antiholomorphic involution exactly when $q_2\in\mathbb R\setminus\{0\}$.
\end{theorem}
\noindent\emph{Uses:} FAM B.4.1, B.4.2 and A.5.1; ARI C.2.1.\par

\begin{proof}
\par\noindent\textit{Step 1.} We check the additional reflection argument, rather than infer a real
structure from the reality of one elliptic invariant. On the period
cover, reflection acts by
\[
 (\tau,\mu,\beta)\longmapsto(-\bar\tau,1-\bar\mu,-\bar\beta).
\]
The theta product and the cusp normalization prove this equality of
continued period functions. For the integral monodromy matrices
$A_1,A_2$ of the retained construction, the lattice comparison is
\[
 H=\begin{pmatrix}1&0&0&0\\0&1&0&0\\-6&0&-1&0\\0&-1&0&-1\end{pmatrix},
 \quad H^2=I,\quad HA_2H=A_2^{-1},\quad
 HA_1H=A_2^{-1}A_1^{-1}A_2,
\]
\[
 \Pi_{-\bar c}(\iota z)H=-\overline{\Pi_c(z)}.
\]
Thus $(z,\zeta)\mapsto(\iota z,-\bar\zeta)$ is an actual torus-family
map. Write $d=(0,0,0,1)^t$, $\varepsilon=(1,2,-4,0)^t$ and
$\varepsilon'=(1,3,-3,0)^t$. Reversing the finite deck generators
changes the target affine vectors to
\[
 v'_1=-\varepsilon+d,\qquad v'_2=\varepsilon'-3d.
\]
The second equality includes the sign of the inverse affine action.
The file \path{Cspecial_reflection.py} checks these matrix identities
over the integers.

\par\noindent\textit{Step 2.} Normalize this reflected datum to the retained vectors
$(\varepsilon,-\varepsilon')$. The linear fibre map is $-I_2$, the
parameter shift is $1$, and the lattice map is $-I+E$,
$E=d(1,0,0,0)$. The differences between target and transported affine
vectors are $2d$ and $-4d$. They lie in the images of the respective
norm operators. Explicit integral correcting cocycles are
$\ell_1=(0,0,-2,0)^t$ and $\ell_2=(0,2,0,0)^t$; the norm equations
are $N_1\ell_1=-2d$ and $N_2\ell_2=4d$. Rational solutions of
$(I-A_j)r_j=(v_j-(-I+E)v'_j)/m_j+\ell_j$ can be chosen with first
coordinate zero: $r_1=(0,-2/3,-2/3,0)^t$ and
$r_2=(0,1,1,0)^t$. These give translations holomorphic across both
root discs. Their cusp cocycle is
$\ell_0=-A_0(\ell_1+A_1\ell_2)=(0,0,4,-2)^t$; its coordinates are
integral, so the corresponding shear and inversion preserve the
periodic triangular cusp fan. They therefore extend across its toric
charts. The differences of the local corrections form a cocycle in
the rank-two vector bundle whose invariant line and quotient are
$\mathcal O$ and $\mathcal O(-1)$ on the base sphere. Both have zero
first cohomology, so the local maps glue. Equivalently one splits the
Laurent series of the cocycle into its nonnegative and negative
parts in these two line-bundle frames. This checks the precise
gluing step of FAM B.4.2 (Statement~\ref{fam:B.4.2}), including the finite affine translations
and cusp extension. The resulting map has target $X_{1-\bar c}$.

\par\noindent\textit{Step 3.} Every other antiholomorphic comparison composed with its inverse is
holomorphic, so Statement~\ref{fam:A.5.1} gives exactly the congruence in the statement.
A self-map exists iff $\operatorname{Re}c\in1/2+\mathbb Z$, iff
$w=-\bar w$, iff $q_2=iw/216$ is real. Its square is a holomorphic
translation $d_0\in\mathbb C/\mathbb Z$, since the automorphism group
is the $\mathbb C^*$ of \cite[theorem \texttt{p2:thm:full-automorphisms}]{DGKXv1}. The constructed reflection acts on
this additive translation coordinate by conjugation. Commutation
with its square gives $d_0=\bar d_0$ in $\mathbb C/\mathbb Z$, hence
$d_0$ has a real representative. Composing the reflection with
translation by $-d_0/2$ makes its square the identity. Necessity and
sufficiency for an involution follow.
\end{proof}

\subsection{Announced real fixed loci}
\begin{definition}[MAPS A.3.setup]\label{maps:A.3.setup}
Here $d=\widehat\delta$, and $\rho_{\rm circ}(b)$ is translation by $be$ for $b\in\mathbb R/\mathbb Z$.
We write $D_1,D_2,D_{12}$ for the three cusp double curves belonging to
edges in directions $e_1,e_2,e_1-e_2$, respectively. Thus $D_2$ is the
fixed sphere of the circle action $\rho_{\rm circ}$, whereas $D_1$ is the double
curve along which the literal character has a limit. This distinction
is essential in the real calculation as well.
\end{definition}
\noindent\emph{Uses:} FAM B.4.1 and the cusp fan.\par

\begin{announcement}[MAPS A.3.1; A]\label{maps:A.3.1}
For a real member the antiholomorphic involutions have exactly two
conjugacy classes under holomorphic automorphisms. For any fixed
involutive splitting $\sigma$ they are represented by $\sigma$ and
$\rho_{\rm circ}(1/2)\sigma$. Every nonempty fixed set is a closed connected
orientable totally real three-manifold with the mod-two homology of
$S^3$. The circle $\rho_{\rm circ}$ preserves each such fixed set and acts
semifreely on it.
\end{announcement}
\noindent\emph{Uses:} FAM B.4.1 and B.4.2; the semifree circle action.\par

The mechanism combines the sign of the coefficient of an involution in the split group $\mathbb C^*\rtimes\mathbb Z/2$ with the fixed-set calculation for the circle action. The statement concerns the smooth fixed manifolds and the two conjugacy classes; the embedded knot type is not determined.

\noindent\emph{Proof:} proof in the repository, MAPS A.3.1 (\url{https://github.com/KokunoYumeto/s6-complex-structure-workbench}).

\begin{definition}[MAPS A.3.representatives]\label{maps:A.3.representatives}
For the explicit representatives, put
\[
 L_F=(2\pi i)^{-1}\log F,\qquad L_G=(2\pi i)^{-1}\log G,
 \qquad F^3=t,\quad G^4=t-1.
 \tag{MAPS A.3.1}
\]
The roots are positive over $t>1$. We continue from that interval
through the upper half $t$-plane. Their logarithms are taken on the
orbifold uniformizing cover. Write $\iota$ for its reflection lift.
Every point of $-i\R_{>0}$ has a representative $c=1/2+iy$, and every
point of $i\R_{>0}$ has a representative $c=-1/2+iy$, with $y$ in the
admissible range. Set
\[
\begin{array}{c|c|c}
 &\alpha&S(t)\\\hline
 c=1/2+iy&1/3&2L_F-4L_G-1/6\\
 c=-1/2+iy&1/2&-2L_G+1/4
\end{array}
\qquad b_c(z)=\alpha\Pi_c(z)\widehat\gamma+S(t(z))e.
\tag{MAPS A.3.2}
\]
Here the fibre vector $e=\Pi_c d=(0,1)^{\mathsf t}$ makes explicit the fibre-coordinate interpretation of $d$ in the source translation formula.
\end{definition}
\noindent\emph{Uses:} FAM B.4.2 and the specified root branches.\par

\begin{announcement}[MAPS A.3.2; A]\label{maps:A.3.2}
On the ordinary period charts the formulas
\[
 \sigma_c[z,\zeta]=[\iota z,\bar\zeta+b_c(\iota z)],
 \qquad \sigma_c^- =\rho_{\rm circ}(1/2)\sigma_c
 \label{maps:eq:A.3.3}
\]
extend to antiholomorphic involutions of the compact threefold.
With these representatives their complete fixed-set data are as follows.
Here $2T^2$ means the disjoint union of two real two-tori; every displayed
circle in a double curve avoids its two triple-point endpoints.
\begin{center}\small
\begin{tabular}{cclccc}
\toprule
$w$&involution&image in $\mathbb{RP}^1$&on $S_1$&on $S_2$&on $W$\\
\midrule
$-ir$, $r>0$&$\sigma_c$&all $\mathbb{RP}^1$&$2T^2$&$2T^2$&$T^2\sqcup S^1_{D_2}\sqcup S^1_{D_{12}}$\\
$-ir$, $r>0$&$\sigma_c^-$&$[1,\infty]$&$\varnothing$&$T^2$&$S^1_{D_1}\sqcup S^1_{D_2}$\\
$ir$, $r>0$&$\sigma_c$&$\varnothing$&$\varnothing$&$\varnothing$&$\varnothing$\\
$ir$, $r>0$&$\sigma_c^-$&$[\infty,1]$ through $0$&$2T^2$&$T^2$&$S^1_{D_1}\sqcup S^1_{D_{12}}$\\
\bottomrule
\end{tabular}
\end{center}
The first interval is the closure of $t>1$; the second is the closure
of $t<1$. On every open real arc contained in the indicated image the
real regular fibre is $2T^2$. No triple point is real for these structures.
\end{announcement}
\noindent\emph{Uses:} MAPS A.3.1 and its explicit representative data; FAM B.4.2.\par

The mechanism is the exact affine fixed-point congruence on the three real base arcs, followed by the finite-root quotient calculation at $0$ and $1$ and the sign calculation for the toric units at the cusp. These are calculations for the stated representatives, including their half-circle translates.

\noindent\emph{Proof:} proof in the repository, MAPS A.3.2 (\url{https://github.com/KokunoYumeto/s6-complex-structure-workbench}).

\begin{announcement}[MAPS A.3.3; A]\label{maps:A.3.3}
Every nonempty fixed set in Theorem \ref{maps:A.3.2} is diffeomorphic to $S^3$.
Its fundamental group is trivial and its integral homology is
$(\Z,0,0,\Z)$. For negative imaginary $w$, the restricted circle action
is the standard semifree action on $S^3$ with fixed circle
$D_2^{\sigma_c}$. For positive imaginary $w$, the nonempty real form
has a free circle action isomorphic, up to circle and orientation
conventions, to the Hopf action.
\end{announcement}
\noindent\emph{Uses:} MAPS A.3.1 and A.3.2.\par

For negative imaginary $w$, the mechanism uses the surface orbit space of the semifree action. For positive imaginary $w$, it uses the two solid-torus caps and the determinant-one comparison of their primitive meridians. No ambient-isotopy classification or Haefliger knot invariant of the embedded three-spheres is asserted.

\noindent\emph{Proof:} proof in the repository, MAPS A.3.3 (\url{https://github.com/KokunoYumeto/s6-complex-structure-workbench}).

\section{The monodromy group and Klingler's question}
\label{sec:monodromy}
The indefinite alternating form and the question about proper discontinuity
are those of \cite[Remark~8.4]{E26}. We retain the actual integral
monodromy of the constructed family.
\begin{definition}[FAM F.setup]\label{fam:F.setup}
We use the clockwise root deck generators and the lattice basis specified
in our conventions.  In this section a group acts properly discontinuously when, for every
compact subset $K$ of its space, only finitely many group elements $g$ satisfy
$gK\cap K\ne\varnothing$.  This definition permits finite stabilizers.  We
write $d=\widehat\delta$, $E=d\otimes\gamma$, and
\[
 J_2=\begin{pmatrix}0&1\\-1&0\end{pmatrix},\qquad
 \Omega=\begin{pmatrix}
 0&0&0&6\\0&0&1&0\\0&-1&0&0\\-6&0&0&0
 \end{pmatrix}.
\]
Thus the alternating form is
$\psi=u_{\Lambda}\wedge w_{\Lambda}+6\gamma\wedge\delta$, an integral alternating form of
type $(1,6)$ on $\Lambda$.  We use $\Omega$ for its matrix and retain $J$ for
a real complex structure.  The domain $\mathcal D$ consists of the real
endomorphisms $J$ of $\Lambda_{\mathbb R}$ satisfying
$J^2=-I$, $J^{\mathsf t}\Omega J=\Omega$, and
$\operatorname{sign}(\Omega J)=(2,2)$.  It is the stated
$\operatorname{Sp}_4(\mathbb R)/U(1,1)$, with its complex structure as an
open subset of the Lagrangian Grassmannian.  An integral symplectic matrix
acts by $J\mapsto gJg^{-1}$.
\end{definition}
\noindent\emph{Uses:} The retained integral marking and monodromy matrices.\par

\subsection{An exact arithmetic description}
\begin{theorem}[FAM F.1.1; T]\label{fam:F.1.1}
Let $P_{d,\mathbb Z}$ be the integral symplectic group fixing $d$ as a
vector, and let $P_{\mathbb Z}$ be the integral symplectic group preserving
$\mathbb Zd\subset\ker\gamma\subset\Lambda$.  Every element of
$P_{d,\mathbb Z}$ has a unique expression
\begin{equation}\label{fam:eq:F-g}
 g(B,x,t)=
 \begin{pmatrix}
 1&0&0\\6x&B&0\\t&-x^{\mathsf t}J_2B&1
 \end{pmatrix},\qquad
 B\in\operatorname{SL}_2(\mathbb Z),\quad x\in\mathbb Z^2,\quad t\in\mathbb Z,
\end{equation}
where the block sizes are $1,2,1$.  Put
\[
 S=\begin{pmatrix}0&-1\\1&0\end{pmatrix},\quad
 T=\begin{pmatrix}1&1\\0&1\end{pmatrix},\quad
 a=\begin{pmatrix}0&1\\-1&-1\end{pmatrix}.
\]
There is a homomorphism
$\alpha:\operatorname{SL}_2(\mathbb Z)\to\mathbb Z/12$ determined by
$\alpha(S)=9$ and $\alpha(T)=1$; in particular $\alpha(a)=4$.
With $Q(x)=x_1+x_2+x_1x_2$, the actual monodromy group is
\begin{equation}\label{fam:eq:F-Gamma}
 \Gamma=\left\{g(B,x,t):t+6Q(x)-\alpha(B)=0\ \text{in }\mathbb Z/12\right\}.
\end{equation}
Consequently $[P_{d,\mathbb Z}:\Gamma]=12$ and
$P_{\mathbb Z}=\{\pm I\}\times P_{d,\mathbb Z}$, so
$[P_{\mathbb Z}:\Gamma]=24$.  Its image in the Levi quotient is exactly
$\operatorname{SL}_2(\mathbb Z)$, with scalar on $d$ equal to $1$.
\end{theorem}
\noindent\emph{Uses:} The matrices $A_1,A_2$ and the integral symplectic form.\par

\begin{proof}
\par\noindent\textit{Step 1.} The matrices in the fixed conventions, in the fixed basis, are
\[
 A_1=\begin{pmatrix}1&0&0&0\\6&0&1&0\\-6&-1&-1&0\\-2&1&0&1\end{pmatrix},
 \qquad
 A_2=\begin{pmatrix}1&0&0&0\\0&0&-1&0\\-6&1&0&0\\3&0&1&1\end{pmatrix}.
\]
Multiplication gives $A_1^3=A_2^4=I$ and
$A_j^{\mathsf t}\Omega A_j=\Omega$.  In particular
$A_1=g(a,(1,-1)^{\mathsf t},-2)$ and
$A_2=g(S,(0,-1)^{\mathsf t},3)$.  Their images $a,S$ generate
$\operatorname{SL}_2(\mathbb Z)$ because $Sa=T$ and $S,T$ perform the
Euclidean algorithm on primitive integral columns.

\par\noindent\textit{Step 2.} A matrix fixing $d$ and preserving $\Omega$ preserves
$d^{\perp}=\ker\gamma$.  Its first row is $(1,0,0,0)$, its last column is
$d$, its middle block $B$ is symplectic on
$d^{\perp}/\mathbb Rd$, and its remaining middle column, say $v$, satisfies
$6r=-v^{\mathsf t}J_2B$.  Integrality of $r$ and unimodularity of $B$
give $v=6x$ with $x\in\mathbb Z^2$.  There is no further condition on
the lower left entry $t$.  This proves \eqref{fam:eq:F-g}.  Direct multiplication
gives
\begin{equation}\label{fam:eq:F-law}
 g(B,x,t)g(C,y,s)
 =g(BC,x+By,t+s-6x^{\mathsf t}J_2By).
\end{equation}

\par\noindent\textit{Step 3.} For completeness, the presentation of the Levi group needed here follows
from an elementary action on the real projective line.  The images of $a$
and $S$ in $\operatorname{PSL}_2(\mathbb Z)$ have orders $3$ and $2$.
On positive real numbers, $a(z)=-1/(z+1)$ has image $(-1,0)$ and
$a^2(z)=-(z+1)/z$ has image $(-\infty,-1)$; $S$ carries negative real
numbers to positive real numbers.  A cyclically reduced alternating word
starting with a nonidentity power of $a$ and ending with $S$ therefore
carries $(-\infty,0)$ into one of these two proper subintervals.  It cannot
be the identity.  Every reduced word is conjugate to such a word or to a
nonidentity element of one factor.  Together with the Euclidean algorithm,
this proves the presentation
$\operatorname{PSL}_2(\mathbb Z)=\langle\bar a,\bar S\mid
\bar a^3=\bar S^2=1\rangle$.  It follows that
\begin{equation}\label{fam:eq:F-SLpresentation}
 \operatorname{SL}_2(\mathbb Z)=
 \langle a,S\mid a^3=S^4=1,\ [a,S^2]=1\rangle.
\end{equation}
Indeed the group on the right modulo its central element $S^2$ is the
just-proved free product, and $S^2$ maps to the nonidentity matrix $-I$.
Its kernel in the displayed matrix group is consequently trivial.
The values $4,9$ respect all these relations, proving the assertion about
$\alpha$; $T=Sa$ has value $1$.

\par\noindent\textit{Step 4.} The quadratic function $Q$ modulo $2$ is zero at $(0,0)$ and one at all
three other points of $\mathbb F_2^2$.  Hence $Q(Bx)=Q(x)\pmod2$ for every
$B\in\operatorname{SL}_2(\mathbb Z)$.  Also
\[
 Q(x+y)-Q(x)-Q(y)\equiv x^{\mathsf t}J_2y\pmod2.
\]
Equation \eqref{fam:eq:F-law} now shows that
$h(g(B,x,t))=t+6Q(x)\pmod{12}$ is a homomorphism.
The homomorphism $\chi=h-\alpha\circ B$ is surjective, since
$\chi(g(I,0,1))=1$.  Its values on $A_1,A_2$ are zero.

\par\noindent\textit{Step 5.} We verify equality with its kernel; containment alone would not determine
the group.  Let $N_{\mathbb Z}$ be the kernel of the map to the middle
block.  The kernel of
$\langle g_1,g_2\mid g_1^3=g_2^4=1\rangle\to
\operatorname{SL}_2(\mathbb Z)$ is normally generated by
$[g_1,g_2^2]$, by \eqref{fam:eq:F-SLpresentation}.  Its image in $\Gamma$ is
therefore normally generated by
\[
 K=[A_1,A_2^2]=g(I,(0,-1)^{\mathsf t},6).
\]
Put
\[
 U=A_2KA_2^{-1}=g(I,(1,0)^{\mathsf t},-6),\quad
 V=K^{-1}=g(I,(0,1)^{\mathsf t},-6),\quad Z=I+12E.
\]
Equation \eqref{fam:eq:F-law} gives $[U,V]=Z^{-1}$.  Thus $U,V,Z$ generate
\begin{equation}\label{fam:eq:F-H}
 H=\{g(I,x,t):t\in6\mathbb Z,\quad t/6\equiv Q(x)\pmod2\}.
\end{equation}
More explicitly,
$U^rV^sZ^k=g(I,(r,s)^{\mathsf t},-6(r+s+rs)+12k)$, so these expressions
give every member of the set.  Conjugation by $g(B,x,t)$ sends
$g(I,y,s)$ to $g(I,By,s-12x^{\mathsf t}J_2By)$; the invariance of $Q$
modulo $2$ proves that $H$ is normal in $P_{d,\mathbb Z}$.  The normal
closure of $K$ is contained in $H$ and contains $U,V,Z$, so it is exactly
$H$.  We have proved $\Gamma\cap N_{\mathbb Z}=H$.  Both $\Gamma$ and
$\ker\chi$ surject onto the same Levi group and have intersection $H$
with $N_{\mathbb Z}$.  Given $g\in\ker\chi$, choose $\gamma\in\Gamma$
with the same middle block.  Then $g\gamma^{-1}\in H$, and $g\in\Gamma$.
This proves \eqref{fam:eq:F-Gamma} and the index $12$.

\par\noindent\textit{Step 6.} Finally an integral automorphism preserving the primitive line
$\mathbb Zd$ sends $d$ to $d$ or $-d$.  Multiplication by the central
matrix $-I$ converts the latter case to the former.  The intersection of
$\{\pm I\}$ with $P_{d,\mathbb Z}$ is trivial.  This proves the direct
product and the index $24$.
\end{proof}

\begin{proposition}[FAM F.1.2; T]\label{fam:F.1.2}
The group in \eqref{fam:eq:F-Gamma} has center
$\langle Z\rangle\cong\mathbb Z$ and is the following explicit central
extension of $\mathbb Z^2\rtimes\operatorname{SL}_2(\mathbb Z)$ by
$\mathbb Z$.  Choose the representative $a(B)\in\{0,\ldots,11\}$ of
$\alpha(B)$.  Its underlying set consists of triples $(B,x,k)$, with
product
\[
 (B,x,k)(C,y,l)=(BC,x+By,k+l+\kappa((B,x),(C,y))),
\]
where
\begin{align*}
12\kappa((B,x),(C,y))={}&a(B)+a(C)-a(BC)-6Q(x)-6Q(y)\\
 &-6x^{\mathsf t}J_2By+6Q(x+By).
\end{align*}
Its real and rational Zariski closures are respectively the full real
and rational symplectic stabilizers of the vector $d$.  They have
dimension $6$, with three-dimensional Heisenberg unipotent radical and
Levi quotient $\operatorname{SL}_2$.
\end{proposition}
\noindent\emph{Uses:} FAM F.1.1 and its exact Heisenberg kernel.\par

\begin{proof}
\par\noindent\textit{Step 1.} The numerator defining $\kappa$ is divisible by $12$, by the two
congruences proved above.  The bijection
\[
 (B,x,k)\longmapsto g(B,x,a(B)-6Q(x)+12k)
\]
intertwines the displayed multiplication with matrix multiplication by
\eqref{fam:eq:F-law}.  This proves associativity as well as the claimed abstract
description, without leaving an unspecified extension class.  The kernel
of the homomorphism to $\mathbb Z^2\rtimes\operatorname{SL}_2(\mathbb Z)$
is $\langle Z\rangle$.  That affine group has trivial center: commuting
with all translations forces its linear part to be $I$, and commuting
with $S,T$ forces its translation part to vanish.  Thus the center of
$\Gamma$ is exactly the stated kernel.

\par\noindent\textit{Step 2.} The logarithms $U-I,V-I,Z-I$ are linearly independent elements of the
three-dimensional Heisenberg Lie algebra.  The square of each is zero.
A polynomial vanishing on all integral powers $I+n(U-I)$ vanishes on
the whole one-parameter line; the same holds for $V$ and $Z$.  Hence the
Zariski closure contains the full unipotent radical.  Its Levi image
contains $S,T$; the powers of $T$ and of $STS^{-1}$ similarly supply the
two opposite unipotent subgroups of $\operatorname{SL}_2$.  Their Lie
algebras and their bracket span $\mathfrak{sl}_2$.  They generate the
connected algebraic group $\operatorname{SL}_2$.  A subgroup containing
the entire unipotent kernel and surjecting onto this Levi quotient is
the entire vector stabilizer.  Its dimension is $3+3=6$.  This reasoning
uses equations over $\mathbb Q$, so proves the rational assertion as
well as the real assertion.
\end{proof}

\subsection{Properness and the maximal open set}
\begin{definition}[FAM F.2.setup]\label{fam:F.2.setup}
Define the real-analytic function
\[
 q_d:\mathcal D\longrightarrow\mathbb R,\qquad q_d(J)=\psi(d,Jd),
 \qquad\mathcal O=\{J:q_d(J)\ne0\}.
\]
The function $q_d$ is invariant under $P_{\mathbb Z}$: replacing $d$ by
$-d$ changes both of its arguments.  Write $\mathcal O_+$ and
$\mathcal O_-$ for its two sign regions.
\end{definition}
\noindent\emph{Uses:} FAM F.setup.\par

\begin{proposition}[FAM F.2.1; T]\label{fam:F.2.1}
The assertion that $\Gamma$ acts properly discontinuously on all of
$\mathcal D$ is false.  The following point of $\mathcal D$ has an
explicit infinite stabilizer:
\[
 J_0=\begin{pmatrix}
 0&0&-1&0\\0&0&0&-6\\1&0&0&0\\0&1/6&0&0
 \end{pmatrix},\qquad
 R=\begin{pmatrix}
 1&0&0&0\\0&1&72&0\\0&0&1&0\\-12&0&0&1
 \end{pmatrix}\in\Gamma.
\]
Every power of $R$ fixes $J_0$, and all these powers are distinct.
\end{proposition}
\noindent\emph{Uses:} FAM F.1.1 and the explicit matrices $J_0,R$.\par

\begin{proof}
Direct multiplication gives $J_0^2=-I$ and
$J_0^{\mathsf t}\Omega J_0=\Omega$.  Moreover
\[
 \Omega J_0=\begin{pmatrix}
 0&1&0&0\\1&0&0&0\\0&0&0&6\\0&0&6&0
 \end{pmatrix}
\]
has eigenvalues $1,-1,6,-6$, so $J_0\in\mathcal D$ and $q_d(J_0)=0$.
The matrix $R$ equals $g(T^{72},0,-12)$ and belongs to \eqref{fam:eq:F-Gamma},
because $\alpha(T^{72})=0\pmod{12}$.  Its commutator with $J_0$ is zero.
Writing $R=I+N$ gives $N\ne0$ and $N^2=0$, so $R^m=I+mN$ and the
powers are distinct.  The compact singleton $\{J_0\}$ therefore violates
the stated definition of proper discontinuity.  These identities are
checked both with SymPy and with a separate rational matrix
implementation in \texttt{code/f\_exact.py}.
\end{proof}

\begin{theorem}[FAM F.2.2; T]\label{fam:F.2.2}
The set $\mathcal O$ is the largest $\Gamma$-invariant open subset of
$\mathcal D$ on which the action is properly discontinuous.  It is also
the largest such open subset for $P_{\mathbb Z}$.  Each sign region is
a complex manifold in the analytic topology, and its quotient by either
discrete group is a Hausdorff complex analytic orbifold.
\end{theorem}
\noindent\emph{Uses:} FAM F.1.1, F.1.2 and F.2.1; the homogeneous-space stabilizer calculation.\par

\begin{proof}
\par\noindent\textit{Step 1.} Let $P_{\mathbb R}$ be the real symplectic group preserving the real line
$\mathbb Rd$.  We first prove properness on each sign region.  For a fixed
$J$, the form
$h_J(v,w)=\psi(v,Jw)+i\psi(v,w)$ is a nondegenerate Hermitian form of
signature $(1,1)$ on the complex vector space $(\Lambda_{\mathbb R},J)$,
with the convention conjugate linear in the first entry.  When
$q_d(J)\ne0$, the complex line $\mathbb Cd$ is nondegenerate and its
Hermitian orthogonal complement is a complex line of the opposite sign.
Normalize $d$ by a positive real scalar and choose a unit vector in this
complement.  Comparing these two Hermitian bases for two complex
structures with the same sign gives a real symplectic map preserving
$\mathbb Rd$ and conjugating one complex structure to the other.  Thus
$P_{\mathbb R}$ is transitive on each sign region.

\par\noindent\textit{Step 2.} An element of the stabilizer of $J$ inside $P_{\mathbb R}$ is complex
linear and sends $d$ to $r d$ for a real nonzero $r$.  Preservation of
$q_d(J)$ gives $r^2=1$.  On the complementary complex line it acts by an
element of $U(1)$.  The stabilizer is consequently the compact group
$\{\pm1\}\times U(1)$.  Each sign region is therefore a homogeneous
space of $P_{\mathbb R}$ with compact stabilizer $L$.  Local sections of
$P_{\mathbb R}\to P_{\mathbb R}/L$, and a finite cover of a compact
subset $K$, give a compact subset $C$ of $P_{\mathbb R}$ mapping onto
$K$.  Every element carrying a point of $K$ back into $K$ belongs to
$CLC^{-1}$, which is compact.  This proves properness.  Restricting a
proper action to a closed discrete subgroup gives
the required finite-intersection condition.  Both $\Gamma$ and
$P_{\mathbb Z}$ are closed discrete matrix groups.  The two sign regions
are disjoint invariant open-and-closed subsets of $\mathcal O$, so the
same finite-intersection condition holds on their union.  Proper actions
by discrete groups have finite stabilizers and Hausdorff quotients; their
holomorphic local finite-group quotient charts supply the asserted
analytic orbifolds.

\par\noindent\textit{Step 3.} It remains to prove maximality, which requires more than one bad point.
The real vector stabilizer $P_{d,\mathbb R}$ acts transitively on
$\mathcal D_0=\{J:q_d(J)=0\}$.  Here is a direct comparison map.  In the
Hermitian space just described, $d$ is a nonzero null vector.  Choose
$v$ with $h_J(d,v)=1$, and replace $v$ by
$v-\tfrac12h_J(v,v)d$.  This preserves $h_J(d,v)=1$ and makes $v$ null.
The ordered complex basis $(d,v)$ has Gram matrix
$\bigl(\begin{smallmatrix}0&1\\1&0\end{smallmatrix}\bigr)$.
The complex-linear map between two such bases fixes $d$ and preserves
$h$, hence its imaginary part $\psi$.  As a real map it belongs to
$P_{d,\mathbb R}$ and conjugates the two complex structures.

\par\noindent\textit{Step 4.} The subgroup $P_{d,\mathbb Q}$ is dense in $P_{d,\mathbb R}$.
For the unipotent coordinates in \eqref{fam:eq:F-g} this is density of
$\mathbb Q^3$ in $\mathbb R^3$.  For $\operatorname{SL}_2$ one may
approximate $a,b,c$ by rational numbers in the chart $a\ne0$ and set
$d_{\rm aux}=(1+bc)/a$; translates of that chart cover the group.  Hence
\[
 \{rJ_0r^{-1}:r\in P_{d,\mathbb Q}\}
\]
is dense in $\mathcal D_0$.
Put $R_0=g(T^6,0,-1)$, so $R_0^{12}=R$ and
$(R_0-I)^2=0$.  At $rJ_0r^{-1}$ the nonzero rational nilpotent
$L=r(R_0-I)r^{-1}$ commutes with the complex structure and satisfies
$L^2=0$.  Choose a positive integer $d_{\rm aux}$ clearing its denominators.
The matrix $I+dL$ and its inverse $I-dL$ are integral, symplectic, and
fix $d$; this follows either by conjugating the identity
$(I+s(R_0-I))^{\mathsf t}\Omega(I+s(R_0-I))=\Omega$ or by expanding it.
The twelfth power of $I+dL$ lies in $\Gamma$, because
$P_{d,\mathbb Z}/\Gamma\cong\mathbb Z/12$.  It has infinite order and
fixes $rJ_0r^{-1}$.  Thus points with infinite $\Gamma$-stabilizer are
dense in $\mathcal D_0$.

\par\noindent\textit{Step 5.} Every open neighborhood in $\mathcal D$ of any point of $\mathcal D_0$
therefore contains a point with infinite stabilizer.  An open set on
which the action is proper cannot meet $\mathcal D_0$.  This proves the
largest-open-set assertion for $\Gamma$.  The same points have infinite
$P_{\mathbb Z}$-stabilizer, so it proves that assertion for the full
parabolic as well.
\end{proof}

\begin{proposition}[FAM F.2.3; T]\label{fam:F.2.3}
The positive sign region admits the global period chart
\[
 \mathcal O_+=\left\{
 (\tau,\mu,\beta)\in\mathbb C^3:
 \operatorname{Im}\tau>0,\quad
 \operatorname{Im}\beta-\frac{6(\operatorname{Im}\mu)^2}
 {\operatorname{Im}\tau}<0\right\},
\]
represented by the matrix $\Pi$ in the question.  In this chart,
\begin{equation}\label{fam:eq:F-qD}
 q_d(J)=-\frac6D.
\end{equation}
Every admissible period $\Pi_c(z)$ belongs to $\mathcal O_+$, for every
admissible $c$.  Thus this single maximal proper-action open set contains
all the period images in question.
\end{proposition}
\noindent\emph{Uses:} The global Siegel chart and the signature condition.\par

\begin{proof}
\par\noindent\textit{Step 1.} Let $J\in\mathcal O_+$.  The two vectors $\widehat w,d$ are complex-linearly
independent in $(\Lambda_{\mathbb R},J)$.  Otherwise
$\widehat w=r d+sJd$ for real $r,s$.  Pairing with $d$ gives
$0=sq_d(J)$; hence $s=0$, contradicting the real linear independence
of $\widehat w,d$.  There is consequently a unique complex-linear
isomorphism $(\Lambda_{\mathbb R},J)\to\mathbb C^2$ sending these two
vectors to $(1,0)$ and $(0,1)$.  The corresponding period matrix has
right block $I$.  The equation
$\Pi\Omega^{-1}\Pi^{\mathsf t}=0$ says that its upper left entry is
six times its lower middle entry, giving exactly the displayed matrix.
This construction is the ordinary holomorphic Grassmannian chart,
so its coordinates and inverse are holomorphic.

\par\noindent\textit{Step 2.} In the real symplectic basis
$(\widehat u,\widehat\gamma/6,\widehat w,d)$ the normalized symmetric
period matrix is
$\bigl(\begin{smallmatrix}\tau&\mu\\\mu&\beta/6\end{smallmatrix}\bigr)$.
Its imaginary part has signature $(1,1)$, equivalently
$6(\operatorname{Im}\mu)^2-
\operatorname{Im}\tau\operatorname{Im}\beta>0$.
Solving $\Pi(Jd)=i\Pi d=(0,i)$ in real coefficients gives
\[
 \gamma(Jd)=
 \frac{\operatorname{Im}\tau}
 {\operatorname{Im}\tau\operatorname{Im}\beta-
 6(\operatorname{Im}\mu)^2}.
\]
Since $\psi(d,v)=-6\gamma(v)$, positivity of $q_d$ now gives
$\operatorname{Im}\tau>0$ and $D<0$, and proves \eqref{fam:eq:F-qD}.
Conversely these inequalities make the imaginary period matrix
invertible with the stated signature and give $q_d>0$.  The admissibility
inequality in the question is $D_c<0$ everywhere, which proves the last
assertion by this exact comparison map.
\end{proof}

\subsection{The curve in the proper quotient}
\begin{theorem}[FAM F.3.1; T]\label{fam:F.3.1}
For every admissible $c$, the marked period map on the covering of
$B^\circ$ is a holomorphic immersion into $\mathcal O_+$.  The induced
map
\[
 p_c:B^\circ\longrightarrow\Gamma\backslash\mathcal O_+
\]
is injective and a holomorphic immersion.  It extends over $t=0,1$ to a
proper holomorphic embedding
\[
 \bar p_c:\mathbb C=B\setminus\{\infty\}
 \lhook\joinrel\longrightarrow\Gamma\backslash\mathcal O_+.
\]
The closure of $p_c(B^\circ)$ in this quotient is exactly
$\bar p_c(\mathbb C)$.
\end{theorem}
\noindent\emph{Uses:} FAM F.2.2 and F.2.3; the finite-root and cusp period expansions.\par

\begin{proof}
\par\noindent\textit{Step 1.} For $B=\bigl(\begin{smallmatrix}a&b\\c&d_{\rm aux}\end{smallmatrix}\bigr)$, changing
the marking by $g(B,x,t)$ replaces $\Pi$ by $\Pi g(B,x,t)^{-1}$ and then
renormalizes its right block.  On the quotient by the complex line of
$d$, the two periods are the row $(\tau,1)$.  Right multiplication by
$B^{-1}$ sends it to $(d_{\rm aux}\tau-c,a-b\tau)$.  Hence
\begin{equation}\label{fam:eq:F-tauaction}
 \tau\longmapsto\frac{d_{\rm aux}\tau-c}{a-b\tau}.
\end{equation}
The denominator never vanishes in the upper half plane.  The function
$j(\tau)$ is therefore a holomorphic invariant of this action and
descends to a holomorphic function $j_\Gamma$ on the proper quotient.
Here $j$ is the elliptic invariant, normalized by
$j(E)=1728\,4A^3/(4A^3+27B^2)$ for $y^2=x^3+Ax+B$; changing a lattice
basis preserves the isomorphism of elliptic curves and hence this
function.

\par\noindent\textit{Step 2.} We derive the required value in the present family.  The elliptic curve
whose period is $\tau$ has the equation stipulated in the question,
\[
 y^2=x^3-3t^3(t-1)x+2t^4(t-1)^2.
\]
Its discriminant is $1728t^8(t-1)^3$ and direct cancellation gives
\begin{equation}\label{fam:eq:F-j}
 j(\tau)=1728t\qquad(t\in B^\circ).
\end{equation}
Equivalently this is the defining lift of the elliptic modular map used
for $\tau$ in the period construction.  The computation is reproduced
exactly in \texttt{f\_exact.py}.  In a local base coordinate $t$, the
derivative of the composite $j\circ\tau$ equals $1728$.  Thus the
derivative of the marked period map is nonzero.  The same identity after
passage to the quotient shows both injectivity and injectivity of the
tangent map of $p_c$: $j_\Gamma/1728$ is a holomorphic left inverse.

\par\noindent\textit{Step 3.} At the two elliptic points, the period functions are holomorphic in the
cyclic covering coordinates $s_1,s_2$, with $t=s_1^3$ and
$t-1=s_2^4$ after choosing local coordinates as in the question.  Their
values are nondegenerate because $D_c<0$ there.  The maps from these
discs to $\mathcal O_+$ are equivariant for $A_1,A_2$.  Composing with
the quotient makes them invariant, so they descend to holomorphic maps
from the base discs.  This is the finite-quotient comparison map: a
holomorphic invariant germ on a disc is a convergent series in $s_j^{m_j}$,
which proves descent on local function rings.  The extensions agree
with $p_c$ off their centers and therefore with each other.  Identity
\eqref{fam:eq:F-j} extends by holomorphy, so the resulting $\bar p_c$ still
has holomorphic left inverse $j_\Gamma/1728$.

\par\noindent\textit{Step 4.} The quotient is Hausdorff by \ref{fam:F.2.2}.  For a compact subset $K$ of it,
$\bar p_c^{-1}(K)$ is closed in $\mathbb C$ and lies in the compact set
$j_\Gamma(K)/1728$.  Thus $\bar p_c$ is proper.  A holomorphic section of
a holomorphic map is locally a closed immersion: in local embeddings
into polydiscs, its coordinate along the base is $t$ and the remaining
coordinates are its holomorphic component functions, whose graph
defines the image.  Applying this to the left inverse shows that
$\bar p_c$ is a holomorphic embedding; properness makes its image
closed.  Since $\mathbb C\setminus\{0,1\}$ is dense in $\mathbb C$, its
image has the asserted closure.
\end{proof}

\begin{proposition}[FAM F.3.2; T]\label{fam:F.3.2}
At the two finite punctures the period map has the holomorphic
finite-cover limits
\[
 \Pi_1=
 \begin{pmatrix}2(2-\rho)&\rho&1&0\\
 \beta_c(z_1)&(2-\rho)/3&0&1\end{pmatrix},\qquad
 \Pi_2=
 \begin{pmatrix}3(1-i)&i&1&0\\
 \beta_c(z_2)&(1-i)/2&0&1\end{pmatrix},
 \quad \rho=e^{\pi i/3}.
\]
Their monodromies are the given order-three and order-four matrices.
At the cusp the curve leaves every compact subset of
$\Gamma\backslash\mathcal O_+$.  On the prescribed cusp branch its
Hodge plane in the compact dual tends to
$\mathbb C\gamma+\mathbb Cu$.  The latter plane has complex dimension
two in common with its conjugate and therefore does not represent a
point of $\mathcal D$.
\end{proposition}
\noindent\emph{Uses:} The finite special values and the cusp-normalized period matrix.\par

\begin{proof}
The finite values of $\tau,\mu$ follow by evaluating their transformation
laws at the fixed points: $\tau(z_1)=\rho$,
$\mu(z_1)=(2-\rho)/3$, $\tau(z_2)=i$,
$\mu(z_2)=(1-i)/2$.  Substitution in $\Pi$ gives the two matrices.
The finite covering comparison and its descent were proved in \ref{fam:F.3.1}.
The unipotent cusp transformation is
$\tau\mapsto\tau-1$, $\mu\mapsto\mu$, $\beta\mapsto\beta+1$.
On the cusp branch, $\operatorname{Im}\tau\to+\infty$,
$\mu\to\mu_0$, and $\beta_c+\tau\to b_0+c$, from the holomorphic
cusp expansions.  Consequently
\[
 D_c=-\operatorname{Im}\tau+O(1),\qquad
 q_d(\Pi_c)\longrightarrow0^+,
 \qquad j_\Gamma(p_c(t))=1728t\longrightarrow\infty.
\]
The last equality alone excludes a convergent subnet in any compact
subset of the proper quotient; the first two identities give the
position relative to the maximal proper-action open set.
The rows of $\Pi_c$, regarded as covectors, are
$6\mu\gamma+\tau u+w$ and $\beta_c\gamma+\mu u+\delta$.
Dividing each by $\tau$ gives limits $u$ and $-\gamma$.
Their exterior product has a nonzero limit, so their span converges
in the Grassmannian to the stated plane.  It is spanned by real
covectors, hence equals its conjugate.  A complex structure has
$H^{1,0}\cap H^{0,1}=0$; here the intersection is the entire
two-dimensional plane.  This explicit comparison proves the assertion
about the compact-dual limit.
\end{proof}

\begin{proposition}[FAM F.3.3; T]\label{fam:F.3.3}
The determinant used in the workbench is generically nonzero, with
the stated zero at $z_1$:
\[
 \det\dot Z=6\dot\mu^2-\dot\tau\dot\beta,\qquad
 \det\dot Z(z_1)=0.
\]
The immersion established in \ref{fam:F.3.1} holds at every point of $B^\circ$
without requiring this determinant to be everywhere nonzero.
\end{proposition}
\noindent\emph{Uses:} The derivative transformation laws and the finite special values.\par

\begin{proof}
In the cusp coordinate $s$, the expansion
$Z=s\bigl(\begin{smallmatrix}0&1\\-1&0\end{smallmatrix}\bigr)+C(t_c)$,
$t_c=e^{2\pi is}$, gives
$\det(dZ/ds)=1+O(t_c)$, so the determinant is not the zero function.
At $z_1$ the fixed-point transformation law of $\mu$ gives
$\dot\mu=-\dot\tau/3$.  One may obtain this by differentiating
$\mu\circ g_1=(1-\mu)/\tau$, using
$g_1'(z_1)=-\rho$ and $\rho^2-\rho+1=0$.
The law
$\beta\circ g_1=\beta+2-6(1-\mu)^2/\tau$ then gives
$\dot\beta=2\dot\tau/3$.  Substitution yields
$6\dot\tau^2/9-2\dot\tau^2/3=0$.
The derivative of a map from a curve to the three-dimensional period
chart is a vector of three coordinate derivatives.  Identity
\eqref{fam:eq:F-j} proves that this vector is nonzero over $B^\circ$;
the determinant tests the rank of the associated two-by-two
Kodaira--Spencer matrix.  The maps compared here are explicitly the
coordinate derivative and its determinant, and the displayed computation
states exactly what each supplies.
\end{proof}

\section{Special fibres, curves and surfaces}
\label{sec:fibres}
The N\'eron--Severi equation below is already
\cite[\S9.1, Lemma~9.2]{AC26}; the accompanying independence and
indefiniteness calculations occur in its Lemmas~9.3--9.4. Its
Proposition~9.5 and Corollary~9.8 give the very-general conclusion
$\operatorname{NS}=\mathbb Z\eta$ and algebraic dimension zero.
Those parts are credited to that source, with a derivation retained here
under the binding T status. The explicit quotient loci, examples and
exhaustive descriptions below have their distinct source statements.
\subsection{Elliptic quotients and the torsion locus}
\begin{proposition}[FAM D.1.1; T]\label{fam:D.1.1}
The specified linear projection \(\operatorname{pr}_1:\mathbb C^2\to
\mathbb C\) induces a holomorphic group homomorphism
\(F_t\to E_t=\mathbb C/(\mathbb Z\tau+\mathbb Z)\) exactly on
\(6\mu=A\tau+B\), \(A,B\in\mathbb Z\). It is then surjective with
connected kernel
\(\mathbb C/(\mathbb Z+(A\mu-\beta)\mathbb Z)\).
On the smooth base the two values of \(t\) are precisely
\[
 t_\pm=\frac{64(47\pm27\sqrt3)}{11}.
\]
\end{proposition}
\noindent\emph{Uses:} The fixed period matrix and the exact division-polynomial computation.\par

\begin{proof}
\par\noindent\textit{Step 1.} The images of the four lattice generators under the linear projection
are \(6\mu,\tau,1,0\). This proves descent in both directions.
The integral basis change replacing \(\widehat\gamma\) with
\(\widehat\gamma-A\widehat u-B\widehat w\) changes its period to
\((0,\beta-A\mu)\); the remaining vertical generator is \((0,1)\).
The other two periods map onto the entire target lattice, so the kernel
is connected. Its second modulus has imaginary part
\(A\operatorname{Im}\mu-\operatorname{Im}\beta=-D>0\).

\par\noindent\textit{Step 2.} The elliptic model may be scaled by the explicit isomorphism
\((X,Y)\mapsto(2X,2\sqrt2Y)\) to
\[
 Y^2=X^3-\frac34t^3(t-1)X+\frac14t^4(t-1)^2,
 \qquad P=(0,t^2(t-1)/2).
\]
Its sixth division polynomial at \(P\) is
\[
 \psi_6(P)=\frac9{65536}t^{24}(t-1)^9
 (11t^2-6016t+8192).
\]
For a smooth elliptic fibre, \(\psi_n(P)=0\) is exactly the condition
\(nP=O\), since the multiplication rational map has these and only these
poles in its affine coordinate. The two roots displayed above are simple
and avoid \(0,1\). We computed the same equation independently by the
chord-and-tangent addition algorithm, obtaining
\[
 \begin{aligned}
 x(6P)&=\frac{(89t^2+4736t-4096)Q(t)}
 {1296(11t^2-6016t+8192)^2},\\
 Q(t)&=881t^4+4745344t^3-12079104t^2\\
 &\hspace{2em}+24641536t-16777216.
 \end{aligned}
\]
The division-polynomial recurrence and the addition algorithm are in
separate supplied programs; all polynomial operations use exact rational
arithmetic. Thus the denominator has no unverified numerical root or
cancellation. The decimal values \(1.36510950191\) and
\(545.543981407\) are only convenient approximations to the exact roots.
\end{proof}

\begin{proposition}[FAM D.1.2; T]\label{fam:D.1.2}
The locus on \(B^\circ\) where \(P(t)\) is torsion is countable and dense.
Every fibre over this locus is an abelian surface. The divisor on the
smooth base cut out by \(6kP=O\) has degree \(3k^2-1\), including local
intersection multiplicity, for every positive integer \(k\).
\end{proposition}
\noindent\emph{Uses:} FAM D.1.1; E26 Lemma 2.1 and its height calculation.\par

\begin{proof}
\par\noindent\textit{Step 1.} On a simply connected open subset write uniquely
\(\mu=a_0\tau+b_0\), \(a_0,b_0\) real. The section is torsion precisely
when \((a_0,b_0)\in\mathbb Q^2\). In a local complex parameter \(z\), direct
differentiation gives
\[
 \det\frac{\partial(a_0,b_0)}{\partial(\operatorname{Re}z,
 \operatorname{Im}z)}
 =-\frac{|\mu'(z)-a_0\tau'(z)|^2}{\operatorname{Im}\tau}.
\]
This determinant vanishes on no open set. On an open set of vanishing,
the two real differential equations imply \(da_0=db_0=0\), hence
\(\mu=a_0\tau+b_0\) with constant real coefficients. Analytic continuation
would give this identity on the cover. The cusp law
\((\tau,\mu)\circ g_0=(\tau-1,\mu)\) then forces \(a_0=0\), and the
\(g_2\)-law makes a constant \(\mu\) impossible. Thus every open subset
contains a point where this real map is a local diffeomorphism, and its
image contains rational pairs. This proves density. For each integer
triple \((n,A,B)\), \(n\ne0\), the function \(n\mu-A\tau-B\) is
holomorphic and not identically zero by the same argument. Its zeros are
discrete; a countable union proves countability.

\par\noindent\textit{Step 2.} For \(N\mu=A\tau+B\), \(A,B,N\in\mathbb Z\), \(N>0\), use the linear
coordinate change \(\zeta'_2=\zeta_2-(A/N)\zeta_1\).
The four integral vectors
\[
 N\widehat u-B\widehat\delta,\quad
 N\widehat w+A\widehat\delta,\quad
 N\widehat\gamma-6A\widehat u-6B\widehat w,\quad
 \widehat\delta
\]
span a sublattice of index \(N^3\). Its first two periods lie on
\(\zeta'_2=0\), and its last two lie on \(\zeta_1=0\).
Both intersections are genuine rank-two lattices because \(D<0\).
Their product quotient is therefore a product of elliptic curves, and
inclusion of this sublattice induces a finite surjective holomorphic
homomorphism onto \(F_t\). A positive integral form on the product can
be summed over its finite translation kernel and multiplied by the
index to descend an integral positive form to \(F_t\). The Riemann
polarization criterion proves projectivity.

\par\noindent\textit{Step 3.} For the degree statement we use the height computation actually read in
\cite[Lemma 2.1 and proof]{E26}: \(\langle P,P\rangle=1/6\),
\(\chi(\mathcal O_S)=1\), the component groups at the two reducible fibres
have orders three and two, and the local corrections for \(P\) are
\(4/3\) and \(1/2\). The application to \(6kP\) is as follows.
The corrected divisor class of a section is its class minus \(O\) and
a fibre multiple, plus the unique rational combination of nonidentity
fibre components that makes it perpendicular to those components.
On restriction to the generic fibre this assignment is additive because
the divisor class \([Q]-[O]\) is the group law in \(\operatorname{Pic}^0\).
The perpendicular correction makes its self-intersection the negative of
the height, so the height is quadratic. The element \(6kP\) is in the
identity component at both reducible fibres; its correction terms vanish.
Hence
\[
 6k^2=(6k)^2\langle P,P\rangle
 =2+2((6kP)\mathbin{\cdot}O),
\]
which gives the formula. There are no missing intersection points at the
three special fibres. The computed rational functions for \(6P\) show
that it has finite affine coordinates at \(0,1\), so it is a nonzero point
of each additive identity component; multiplication by \(k\) remains
nonzero in \(\mathbb C\). At infinity its multiplicative coordinate is
\((-a)^{6k}=a^{6k}\ne1\). Thus all intersections counted by the height
occur on \(B^\circ\). This derives the cited height input in precisely
the form needed here; the first five degrees and the correction matrices
are checked exactly in \texttt{cd\_torsion\_independent.py}.
\end{proof}

\noindent The following counterexample addresses only the unrestricted wording in FAM Lead~6. The first-projection assertion of Statement~\ref{fam:D.1.1} stands.\par
\begin{theorem}[FAM D.1.3; T]\label{fam:D.1.3}
The assertion in Lead 6 that \emph{existence of a surjection onto the
specified \(E_t\) with elliptic kernel} is equivalent to \(6P(t)=O\)
is false. It remains true for the specified projection in \ref{fam:D.1.1}.
\end{theorem}
\noindent\emph{Uses:} FAM D.1.1 and the source bound FAM C.2.1.\par

\begin{proof}
\par\noindent\textit{Step 1.} Set \(t=32/5\). The exact fourth division polynomial is
\[
 \psi_4(P)=-\frac1{32}t^{10}(t-1)^4(5t-32),
\]
whereas \(\psi_2(P)=t^2(t-1)\ne0\). Thus \(P\) has exact order four.
On the rectangular branch write \(\tau=iT\), \(T>1\), with
\(j(iT)=1728\cdot32/5\). The chosen zero has real part \(1/2\), negative
imaginary part of absolute value less than \(T/2\), and order four;
therefore \(\mu=(2-\tau)/4\). Choose the concrete parameter
\[
 c=\left(\frac98-\frac{101}{4}\right)\tau-\frac94
       -\beta^{\rm part}(t).
\]
It is globally admissible. Indeed the positive-real Jacobi product at this
\((\mu,\tau)\) gives
\(\operatorname{Im}\beta^{\rm part}\geq-T-(3/\pi)\log3\).
Consequently
\(\operatorname{Im}c\leq-185T/8+(3/\pi)\log3<-22<-M\), by C.2.1.
This inequality proves admissibility without using the numerical estimate
\(c=-2.25-34.02109724218209472308\,i\).

\par\noindent\textit{Step 2.} The holomorphic linear map \(L:\mathbb C^2\to\mathbb C\),
\(L(\zeta)=3\zeta_1+4\zeta_2\), takes the four periods to
\[
 -101\tau,\qquad2\tau+2,\qquad3,\qquad4.
\]
These generate exactly \(\mathbb Z\tau+\mathbb Z\): the last two generate
\(1\), and the first two then generate \(\tau\), since
\(\gcd(101,2)=1\). Thus \(L\) descends to a surjective holomorphic group
homomorphism \(F_{32/5}\to E_{32/5}\) with connected one-dimensional
kernel. The kernel is compact as a closed subset of a compact torus,
hence an elliptic curve. Nevertheless \(6P=2P\ne O\).
Every equality, the division-polynomial noncancellation, and the four
lattice images are checked by exact arithmetic in the supplied code.
\end{proof}

\begin{definition}[FAM D.1.Hermitian]\label{fam:D.1.Hermitian}
The real vector-space isomorphism
$\Pi:\mathbb R^4\to\mathbb C^2$ transfers $E$ to the actual alternating
form $\widetilde E(v,w)=E(\Pi^{-1}v,\Pi^{-1}w)$ on the tangent space.
Equation \eqref{fam:eq:cd-NS} gives $J$-invariance. The explicit comparison is
\[
 H_E(v,w)=\widetilde E(v,Jw)-i\widetilde E(v,w).
\] We use $T=\operatorname{Im}\tau$, $U=\operatorname{Im}\mu$ and $D=\operatorname{Im}\beta-6U^2/T$ in the matrix calculation below.
\end{definition}
\noindent\emph{Uses:} FAM D.1.6, Hermitian-form convention.\par

\begin{proposition}[FAM D.1.4; T]\label{fam:D.1.4}
For every smooth fibre and every admissible \(c\), its Neron--Severi lattice
is the following explicit subgroup of \(\mathbb Z^6\):
\begin{equation}\label{fam:eq:cd-NS}
 \left\{E:\ E_{12}-E_{13}\tau+(6E_{23}-E_{14})\mu+
 E_{24}\beta+E_{34}(6\mu^2-\tau\beta)=0\right\}.
\end{equation}
The complex orientation has intersection form
\(E^2=-2\operatorname{Pf}(E)\), where
\(\operatorname{Pf}(E)=E_{12}E_{34}-E_{13}E_{24}+E_{14}E_{23}\).
In particular the invariant class \(\eta=6\gamma\wedge\delta+u_{\Lambda}\wedge w_{\Lambda}\)
has \(\int_{F_t}\eta^2=-12\). Its Hermitian form is indefinite.
The fibres with lattice exactly \(\mathbb Z\eta\) have no compact curves
and algebraic dimension zero.
\end{proposition}
\noindent\emph{Uses:} The period matrix; the Hermitian-form convention in FAM D.1.6; AC26 Lemmas 9.2 and 9.4.\par

\begin{proof}
\par\noindent\textit{Step 1.} On the universal cover the holomorphic two-form is the wedge of the two
rows of \(\Pi\). Its six coefficients in the integral coframe are
\(6\mu^2-\tau\beta,-\beta,6\mu,-\mu,\tau,1\).
An integral real two-form is of type \((1,1)\) exactly when its wedge
with this form vanishes: its \((0,2)\) part is then zero, and reality
also annihilates its \((2,0)\) part. Expanding the wedge gives
\eqref{fam:eq:cd-NS}. The exponential sequence identifies these integral
\((1,1)\) classes with Chern classes of holomorphic line bundles, proving
the lattice assertion.

\par\noindent\textit{Step 2.} The real period matrix in C.3.1 has determinant \(TD<0\). Thus the
integral orientation \((\widehat\gamma,\widehat u,\widehat w,
\widehat\delta)\) is opposite to the complex orientation. Algebraically
\(E\wedge E=2\operatorname{Pf}(E)\gamma\wedge u_{\Lambda}\wedge w_{\Lambda}\wedge\delta\),
which proves the sign in the displayed intersection formula. In
particular the expression ``\(\eta^2=12\)'' in \cite[Lemma~9.4]{AC26} is correct
as its displayed coefficient of \(\gamma\wedge u_{\Lambda}\wedge w_{\Lambda}\wedge\delta\),
and has the opposite sign as a complex-oriented numerical
self-intersection. We retain the coefficient identity and record the
orientation comparison explicitly.

\par\noindent\textit{Step 3.} For the precise Hermitian convention $H_E$ defined in \ref{fam:D.1.6}, direct
substitution of the four periods gives $H_\eta=-2G^{-1}$, where
\[
 G=\begin{pmatrix}2T&2U\\2U&\operatorname{Im}\beta/3\end{pmatrix},
 \qquad \det G=2TD/3<0.
\]
Thus $H_\eta$ has one positive and one negative eigenvalue, and every
nonzero multiple remains indefinite. The integration
current of an effective divisor, averaged under translations of the
compact torus, is a nonzero semipositive constant \((1,1)\)-form in its
Chern class. It cannot be an indefinite multiple of \(\eta\).
Thus no curve exists when \(\operatorname{NS}=\mathbb Z\eta\).
A nonconstant meromorphic function would have a nonempty polar divisor
(after subtracting a constant if necessary), so the algebraic dimension
is zero.
\end{proof}

\begin{proposition}[FAM D.1.5; T]\label{fam:D.1.5}
For a fixed \(c\), all but countably many smooth fibres have
\(\operatorname{NS}=\mathbb Z\eta\). Conversely, for every fixed smooth
\(t\), there are admissible parameters tending to the cusp of the
parameter disc for which \(F_t\) is an abelian surface, irrespective of
the order of \(P(t)\).
\end{proposition}
\noindent\emph{Uses:} FAM D.1.4; the five-function independence of AC26 Lemma 9.3.\par

\begin{proof}
\par\noindent\textit{Step 1.} We verify the required independence of five holomorphic functions.
Suppose \(A+B\tau+C\mu+D\beta+E(6\mu^2-\tau\beta)=0\).
Apply \(g_0\), using \((\tau,\mu,\beta)\circ g_0=(\tau-1,\mu,\beta+1)\),
and subtract. The result is \(-B+D+E(\beta-\tau+1)=0\).
A second cusp difference gives \(2E=0\), hence \(E=0,D=B\).
Applying \(g_2\) to the remaining relation and multiplying the difference
by \(-\tau\) gives
\[
 6B\mu^2+C(\tau-1)\mu+B(\tau^2+3\tau+1)-C\tau=0.
\]
Its cusp difference gives \(C\mu+2B(\tau+1)-C=0\), whose next cusp
difference forces \(B=0\). Then \(C\mu=C\); the \(g_2\)-law forces
\(C=0\), and finally \(A=D=0\). Thus an integral class outside
\(\mathbb Z\eta\) has a nonidentically-zero holomorphic left side in
\eqref{fam:eq:cd-NS}; its zeros are discrete. Countability of \(\mathbb Z^6\)
proves the first assertion, with explicit verification of its reduction.

\par\noindent\textit{Step 2.} For the second assertion put
\(\beta_n=-3n^2\tau+12n\mu\),
\(c_n=\beta_n-\beta^{\rm part}(t)\). Its imaginary part tends to
minus infinity, so the already proved finite bound on \(M\) makes every
sufficiently large \(n\) admissible. The integral alternating matrix
with nonzero upper entries
\[
 E_{13}=-3n^2,\quad E_{14}=6n,\quad E_{23}=-n,\quad E_{24}=1
\]
satisfies \eqref{fam:eq:cd-NS} exactly at \(\beta_n\), and
\(\operatorname{Pf}(E)=-3n^2<0\). Its complex-oriented square is positive.
A nondegenerate Hermitian form in complex dimension two with positive
determinant is definite. Choosing its sign gives a positive integral
Riemann form, so \(F_t\) is projective. This supplies an explicit family
of projective fibres even at any chosen non-torsion value of the section.
\end{proof}

\begin{theorem}[FAM D.1.6; T]\label{fam:D.1.6}
Fix any admissible $c$. For each integral alternating matrix $E$, let
$Z_E(c)$ be the zero set on the marking cover of the explicit holomorphic
function in \eqref{fam:eq:cd-NS}, with $\beta=\beta^{\rm part}+c$ given by
\eqref{fam:eq:cd-beta}. Let $\pi$ be the cover map to $B^\circ$, and define
\[
 \mathcal E_c=\pi\!\left(\bigcup_{E\ne0,\ \operatorname{Pf}(E)=0}Z_E(c)\right),
 \qquad
 \mathcal A_c=\pi\!\left(\bigcup_{\operatorname{Pf}(E)<0}Z_E(c)\right).
\]
These give the exact loci of all elliptic quotients and all abelian
surfaces, respectively. The algebraic dimension of every smooth fibre is
\[
 a(F_t)=
 \begin{cases}
 2,&t\in\mathcal A_c,\\
 1,&t\in\mathcal E_c\setminus\mathcal A_c,\\
 0,&t\notin\mathcal E_c\cup\mathcal A_c.
 \end{cases}
\]
Together with \eqref{fam:eq:cd-NS}, these are an exact classification for every
$t$, by countable unions of explicitly given analytic zero sets; they do
not assert that the loci have a finite algebraic description. They include
all quotient maps, without requiring preservation of the first coordinate.
\end{theorem}
\noindent\emph{Uses:} FAM D.1.4; the Appell--Humbert and Riemann criteria.\par

\begin{proof}
\par\noindent\textit{Step 1.} We apply the Appell--Humbert correspondence and the Riemann projectivity
criterion in \cite[Theorems 4.1 and 4.4]{Schnell19}; the reduction to the
present explicit periods is as follows. The real vector-space isomorphism
$\Pi:\mathbb R^4\to\mathbb C^2$ transfers $E$ to the actual alternating
form $\widetilde E(v,w)=E(\Pi^{-1}v,\Pi^{-1}w)$ on the tangent space.
Equation \eqref{fam:eq:cd-NS} gives $J$-invariance. The explicit comparison is
\[
 H_E(v,w)=\widetilde E(v,Jw)-i\widetilde E(v,w).
\]
Direct substitution gives $H_E(Jv,w)=iH_E(v,w)$ and
$H_E(w,v)=\overline{H_E(v,w)}$. Its associated real Chern form is
$-\operatorname{Im}H_E=\widetilde E$, since
$(i/2)(H_E(v,w)-H_E(w,v))=\widetilde E(v,w)$.
Thus the Appell--Humbert integrality condition is exactly integrality
of the given $E$, with its sign fully specified. The complex-oriented square computed in \ref{fam:D.1.4} shows that
$\operatorname{Pf}(E)<0$ is equivalent to definiteness of this Hermitian
form. Its sign can be chosen positive. The cited projectivity criterion
then gives precisely the asserted abelian locus, in both directions.

\par\noindent\textit{Step 2.} For nonzero $E$ with Pfaffian zero, its real rank is two. Its nullspace
$K_{\mathbb R}\subset\mathbb R^4$ is a rational plane because it is the
kernel of an integral matrix. $J$-invariance makes $\Pi K_{\mathbb R}$ a
complex line, and $K=K_{\mathbb R}\cap\mathbb Z^4$ is saturated. The maps
\[
 (\Pi K_{\mathbb R})/\Pi K\ \hookrightarrow F_t
 \ \longrightarrow\
 \mathbb C^2/(\Pi K_{\mathbb R}+\Pi\mathbb Z^4)
\]
are holomorphic group morphisms; the first is a closed embedded elliptic
curve and the second is its quotient onto an elliptic curve, with
connected kernel. A basis of the saturated $K$ extends to an integral basis of $\mathbb Z^4$.
The other two periods project to real-linearly independent vectors in the
one-dimensional complex quotient, because $\Pi$ is a real isomorphism.
They therefore form a discrete lattice there. Compactness follows by
taking the image of its compact real fundamental parallelogram. Conversely any elliptic quotient
pulls an integral positive area form back to a nonzero rank-two integral
$(1,1)$ form. This proves the characterization of $\mathcal E_c$.
Equivalently all its kernels can be found by saturated rank-two lattices
with a basis $v,w$ satisfying
$\det_{\mathbb C}(\Pi v,\Pi w)=0$ and real independence. This is a second
exact integer test, with explicit inclusion and quotient morphisms.

\par\noindent\textit{Step 3.} For the algebraic-dimension assertion, an effective divisor on a torus
has a semipositive Hermitian form: averaging its integration current under
translations gives that constant form. A nonzero semipositive form has
rank one or two. Rank two gives the abelian case. Rank one gives the
elliptic quotient just constructed. In the absence of a definite form,
there cannot be two different such kernel subtori. Indeed addition from
their product onto $F_t$ would be a surjective holomorphic homomorphism
with finite kernel, and the product polarization would descend after
multiplication by an integer, giving a definite form. Thus there is a
single kernel $K_t$ in this case. Every effective divisor is a sum of
its translates: its intersection number with $K_t$ is zero, whereas a
curve mapping nonconstantly to the elliptic quotient has strictly positive
intersection with a general quotient fibre. Every meromorphic function
therefore factors through this quotient. One can check the last statement
on a fibre avoiding its polar divisor, where it is a holomorphic function
on a compact elliptic curve and hence constant; analytic continuation
then gives the factorization. The quotient supplies nonconstant
meromorphic functions, so the algebraic dimension is exactly one.
If neither locus occurs, a nonconstant meromorphic function would have
a polar divisor of one of the two forbidden ranks. This proves dimension
zero. Finally monodromy acts by integral determinant-one changes of basis,
preserving \eqref{fam:eq:cd-NS} and its Pfaffian. Thus the stated unions descend
without any marking ambiguity. Each individual matrix and every locus
is specified; no expected identification is left in these formulas.
\end{proof}

\begin{proposition}[FAM D.1.7; T]\label{fam:D.1.7}
There is an explicit nonempty locus of admissible members with ordinary
fibres of algebraic dimension one. On a sufficiently small punctured
cusp neighbourhood set
\[
 c(z)=-6\tau(z)-12\mu(z)-\beta^{\rm part}(z).
\]
Outside a countable set of $z$, the fibre of $X_{c(z)}$ over $\pi(z)$ has
Neron--Severi lattice $\mathbb Z\eta+\mathbb ZE$, where the nonzero upper
entries of $E$ are
$E_{13}=-6$, $E_{14}=-6$, $E_{23}=1$, $E_{24}=1$.
Its algebraic dimension is one. The section $P(\pi(z))$ can simultaneously
be required to be non-torsion.
\end{proposition}
\noindent\emph{Uses:} FAM D.1.6 and the cusp expansions.\par

\begin{proof}
At the specified parameter $\beta=-6\tau-12\mu$. Both $E$ and $\eta$
satisfy \eqref{fam:eq:cd-NS}, and exact multiplication gives
$\operatorname{Pf}(a\eta+bE)=6a^2$.
For any further integral class $F$, substitution in its period equation
gives a linear relation among $1,\tau,\mu,(\mu+\tau)^2$.
These four functions are linearly independent: applying $g_0$ to a
relation $A+B\tau+C\mu+D(\mu+\tau)^2=0$ and subtracting yields
$-B+D(1-2(\mu+\tau))=0$. A further cusp difference forces $D=0$,
and then $B=0$; nonconstancy of $\mu$ forces $A=C=0$.
Comparison of coefficients now gives
$F_{34}=F_{12}=0$, $F_{13}=-6F_{24}$, and
$6F_{23}-F_{14}=12F_{24}$, exactly
$F\in\mathbb Z\eta+\mathbb ZE$. Every other integral class therefore
has a discrete zero set along this period locus. Excluding their
countable union gives the asserted exact Neron--Severi lattice.
It contains the nonzero isotropic $E$ and no matrix of negative
Pfaffian; \ref{fam:D.1.6} gives algebraic dimension one.
Finally
$c=-5\tau-12\mu-(\beta^{\rm part}+\tau)$ has imaginary part tending
to minus infinity at the cusp. A sufficiently small cusp neighbourhood
is therefore admissible by C.2.1. The torsion locus of $P$ is countable
by \ref{fam:D.1.2}, so its exclusion still leaves such points in every open subset
of that neighbourhood. This supplies actual members and fibres, with
all their parameter functions specified.
\end{proof}

\subsection{Compact curves and surfaces}
The cusp normalization and its boundary identifications used here are due to \cite[\S4.6, Proposition~4.6]{AC26}.
\begin{proposition}[FAM D.2.1; T]\label{fam:D.2.1}
For a standard-twist member \(X_c\) with underlying manifold \(S^6\),
every compact irreducible complex curve lies in one fibre of \(f\), and
every compact irreducible complex surface is one of the reduced fibres.
\end{proposition}
\noindent\emph{Uses:} $X_c\simeq S^6$; the K\"ahler geometry of the reduced fibres.\par

\begin{proof}
The given diffeomorphism with \(S^6\) implies \(H^2(X_c;\mathbb Z)=0\).
For a compact irreducible curve \(C\), let
\(\nu:\widetilde C\to C\hookrightarrow X_c\) be its normalization map.
The holomorphic map \(f\nu\) has
\(\deg (f\nu)^*\mathcal O_{\mathbb P^1}(1)=0\), since its Chern class is
pulled back from the zero class on \(X_c\). A nonconstant holomorphic map
of compact Riemann surfaces has positive degree: a regular value has a
positive number of preimages, all with positive local multiplicity.
Therefore \(f\nu\), and hence \(f|_C\), is constant.

An irreducible surface \(Z\) in the smooth threefold is an effective
Cartier divisor. If its image were all of \(B\), a general smooth fibre
would meet it in a nonempty effective divisor \(D\) whose line bundle
is \(\mathcal O_{X_c}(Z)|_{F_t}\). Its integral Chern class is zero.
Choose a flat K\"ahler form \(\omega\) on the torus \(F_t\). The equality
\(\int_{F_t}c_1(D)\wedge\omega=\int_D\omega\) has left side zero and right
side positive, a contradiction. Thus \(Z\) lies in a single fibre.
Each reduced fibre is irreducible of dimension two, so containment forces
equality. This uses the full map into \(X_c\), not a transfer of a
vanishing claim from its normalization.
\end{proof}

\begin{proposition}[FAM D.2.2; T]\label{fam:D.2.2}
There is no rational curve in \(X_c\setminus W_c\). There are elliptic
curves outside all three special fibres, in particular in each of the
two smooth fibres \(F_{t_\pm}\). All curves in \(W_c\) are precisely the
images of irreducible algebraic curves in \(\mathrm{dP}_6\) under its
finite normalization map. Distinct curves have the same image only for
the three paired boundary lines. All curves in a bielliptic fibre
\(S_j\) are images of irreducible curves in its finite product cover.
\end{proposition}
\noindent\emph{Uses:} FAM D.1.1 and D.2.1; the finite normalization and boundary identifications of AC26 Proposition 4.6.\par

\begin{proof}
\par\noindent\textit{Step 1.} By \ref{fam:D.2.1} every curve is vertical. A holomorphic map from
\(\mathbb P^1\) into a torus lifts to its vector-space universal cover,
since \(\mathbb P^1\) is simply connected, and its coordinate functions
are constant by compactness. Thus a smooth torus fibre contains no
rational curve. A map from \(\mathbb P^1\) into a bielliptic surface
lifts through its finite unramified torus cover by the same covering
criterion; the lifted map is constant. These arguments prove the first
assertion, including possibly singular rational images.
\ref{fam:D.1.1} gives an embedded connected elliptic kernel in each \(F_{t_\pm}\),
proving the second.

\par\noindent\textit{Step 2.} For the third assertion let \(\eta:\mathrm{dP}_6\to W_c\) be the finite
normalization morphism. Each component of the inverse image of a compact
curve is compact and one-dimensional and maps onto that curve, because
finite maps have zero-dimensional fibres. It is algebraic in the
projective surface \(\mathrm{dP}_6\). Conversely its image is a compact
analytic curve by properness. Away from the conductor, \(\eta\) is a
biholomorphism, so two distinct curves can have the same image only when
both lie in the conductor; its identifications are exactly the three
opposite pairs computed in FAM C.4.1; see also \cite[\S4.6, Proposition~4.6]{AC26}. This is also an explicit algebraic
parametrization: the curves of \(\mathrm{dP}_6\) are the three exceptional
lines and strict transforms of irreducible plane curves under the blow-up
of the three coordinate points. Their homogeneous equations and their
three multiplicities specify the strict transforms. The finite-product
cover statement for \(S_j\) follows from the identical proper-image and
finite-preimage argument applied to its unramified morphism
\(E_j\times F_j\to S_j\); curves in the product are algebraic.
\end{proof}

\begin{theorem}[FAM D.2.3; T]\label{fam:D.2.3}
The following descriptions give all compact irreducible curves of $X_c$.
In $W_c$ and $S_j$ they are exactly the images described in \ref{fam:D.2.2}.
In a smooth fibre of algebraic dimension zero there are none. In a
smooth fibre of algebraic dimension one they are exactly the translates
of its unique elliptic kernel $K_t$. In an abelian fibre they are exactly
the translates of elliptic subtori, together with the irreducible zero
divisors of the positive Appell--Humbert theta sections described below.
These descriptions, \ref{fam:D.1.6}, and \ref{fam:D.2.1} constitute a classification including
every smooth exceptional fibre and every compact irreducible surface.
\end{theorem}
\noindent\emph{Uses:} FAM D.1.6 and D.2.1; the Appell--Humbert theorem and its theta equation.\par

\begin{proof}
\par\noindent\textit{Step 1.} All curves are vertical by \ref{fam:D.2.1}, so the list of fibre types is exhaustive.
The dimension-zero and dimension-one assertions, including uniqueness,
were proved in \ref{fam:D.1.6}. Let now $F_t=\mathbb C^2/\Pi\Lambda$ be abelian.
For every positive definite Hermitian form $H$ whose imaginary part takes
integer values on $\Pi\Lambda$, and every function
$\alpha:\Pi\Lambda\to U(1)$ satisfying
\[
 \alpha(\lambda+\nu)=\alpha(\lambda)\alpha(\nu)
                 e^{\pi i\operatorname{Im}H(\lambda,\nu)},
\]
take the nonzero entire solutions $\theta:\mathbb C^2\to\mathbb C$ of
\begin{equation}\label{fam:eq:cd-theta-sections}
 \theta(\zeta+\lambda)=\alpha(\lambda)
 e^{\pi H(\zeta,\lambda)+\pi H(\lambda,\lambda)/2}\theta(\zeta).
\end{equation}
The cocycle identity follows by multiplying the displayed factors; its
integral imaginary cross-term has exponential one. It defines a
holomorphic line bundle by quotienting $\mathbb C^2\times\mathbb C$.
The invariant zero set of $\theta$ descends to a compact divisor on
$F_t$. Retain precisely those zero divisors that are reduced and
irreducible. This is an exact parametrization, with all periods and
multiplier equations given; multiplying $\theta$ by a nonzero scalar
does not change its divisor.

\par\noindent\textit{Step 2.} The correspondence \cite[Theorem 4.1 and equation (4.3)]{Schnell19} proves
that every holomorphic line bundle and every section on this actual
quotient has this form. For the cohomology matrix the exact comparison is
$H_E=\widetilde E(\,\cdot\,,J\,\cdot\,)-i\widetilde E$ from \ref{fam:D.1.6};
we choose the sign of $E$ making $H_E$ positive. Its imaginary part is
$-\widetilde E$, which gives the integral semicharacter equation displayed
above. A divisor class with positive Hermitian rank two gives exactly
these positive theta sections. A divisor of Hermitian rank one is pulled
back from its elliptic quotient and, when irreducible and reduced, is a
single fibre. Thus it is a translate of an elliptic subtorus. The null
lattice is computed by the explicit integer kernel in \ref{fam:D.1.6}. No other
rank is possible for an effective nonempty divisor, by positivity of its
averaged current. This proves exhaustiveness on the abelian fibre.

\par\noindent\textit{Step 3.} The same equations give the curves in the product covers of $S_j$:
compute them on its two-elliptic-curve lattice and take images under the
specified finite morphism. Conversely each curve downstairs is the image
of a component of its finite inverse image, so that this procedure omits
none. In $W_c$ the analogous procedure uses the homogeneous polynomial
parametrization of the normalization in \ref{fam:D.2.2}, with the three boundary
pairings already identified. The maps used throughout are finite or
closed embeddings, so they preserve compactness and give one-dimensional
images; no assertion is transferred from a normalization in an invalid
direction.
\end{proof}

\section{The outer boundary and the untwisted filling conjecture}
\label{sec:ends}
\subsection{The intrinsic obstruction at the outer circle}
Throughout this subsection the twists are the standard ones.
Since $q_2=iw/216$, its boundary circle is $|q_2|=1$.
\begin{definition}[PUSH D.1.setup]\label{push:D.1.setup}
We use the standard twists throughout this section. Put
\(H_*=(3/\pi)\log6\), fix a boundary parameter
\(c_*=x+iH_*\), and set \(a=x+3/2\).
\end{definition}
\noindent\emph{Uses:} The exact radius 216.\par

\begin{proposition}[PUSH D.1.1; T]\label{push:D.1.1}
For \(0<|w|<216\), the reduced multiplicity-four surface has
\[
 j(\operatorname{Alb}S_2(w))=J(iw/216),\qquad
 J(q)=q^{-1}+744+196884q+\cdots.
\]
This holomorphic function has no meromorphic continuation through any
point of \(|w|=216\).
\end{proposition}
\noindent\emph{Uses:} FAM C.1.1 and A.4.1; ARI C.2.1; modular invariance of j.\par

\begin{proof}
The fixed-coordinate periods on the product cover, including the
isogeny kernel and the affine deck translation, generate
\(\frac12\mathbb Z+\frac{\kappa_2}{4}\mathbb Z\).
Multiplication by two therefore identifies the Albanese lattice with
\(\mathbb Z+(\kappa_2/2)\mathbb Z\). This verifies the precise step of
FAM A.4.1 used here; taking only the invariant eigenspace would omit the
half-period. The values in FAM C.1.1 and \(\mu(z_2)=(1-i)/2\) give
\[
 \kappa_2=c+\frac32-iH_*,\qquad
 e^{-\pi i\kappa_2}=e^{-3\pi i/2}e^{-\pi H_*}e^{-\pi ic}
 =\frac{iw}{216}.
\]
The positive-imaginary representative is \(-\kappa_2/2\), proving the
identity with the stated Fourier normalization.

For completeness we reproduce the obstruction in FAM C.3.4, rather than regard a numerical Fourier plot as evidence. At a
reduced rational \(r=b/d\) choose a matrix in
\(\operatorname{SL}_2(\mathbb Z)\) with bottom row \((d,-b)\).
For \(z=r+iy\), its transformed imaginary part is \(1/(d^2y)\).
Modular invariance and the convergent Fourier series imply
\(\log|j(r+iy)|=2\pi/(d^2y)+o(1)\).
A meromorphic function near \(r\) grows at most polynomially in
\(1/y\), so no continuation exists there. Every real interval contains
such a rational, excluding continuation at irrational points as well.
A local logarithm of \(q=iw/216\) transfers this result to the circle.
This is a derivation of the earlier natural-boundary result, with the
normalizing constant checked exactly.
\end{proof}

\begin{lemma}[PUSH D.1.2; T]\label{push:D.1.2}
Let \(V\) be a connected Riemann surface and
\(\pi:\mathcal X\to V\) a flat proper holomorphic map of reduced
complex spaces. Suppose that over a nonempty open subset \(V_0\), the
fibres are smooth compact threefolds whose unique anticanonical divisor
is \(2S_v\), where \(S_v\) is a smooth connected surface with
\(b_1(S_v)=2\). Then there is an open subset \(V^{\rm good}\subset V\), with at
most countable complement, on which the reduced support of the
relative anticanonical divisor is a smooth proper surface family. Its
Albanese elliptic invariant is holomorphic.
\end{lemma}
\noindent\emph{Uses:} Grauert direct-image and proper-mapping theorems with the 1963 correction; Ehresmann; Buchdahl Theorem 11.\par

\begin{proof}
\par\noindent\textit{Step 1.} Flatness and the smoothness of the fibres over \(V_0\) imply that
\(\pi\) is a submersion there: the local analytic smoothness criterion
applies because the base is smooth. The relative nonsmooth locus is a
closed analytic subset, by the Jacobian criterion. Its proper image is
analytic by \cite[\S7, Satz~1]{Grauert1960D}, the proper mapping theorem.
It misses \(V_0\), hence is discrete. We remove it and work over the
connected remaining curve \(V'\). Removing a discrete subset from a
connected Riemann surface does not disconnect it. Thus
\(\mathcal X'\to V'\) is a smooth proper map.

\par\noindent\textit{Step 2.} Write \(L=K_{\mathcal X'/V'}^{-1}\). Grauert's direct-image theorem \cite[\S6, Hauptsatz~I]{Grauert1960D} makes \(F=\pi_*L\) and \(R^1\pi_*L\)
coherent. We prove the needed base change directly over the smooth
curve; we do not use the originally stated \S7, Satz~2, whose claimed
closed exceptional locus was corrected in Grauert's 1963 erratum \cite[item~2, p.~36]{Grauert1963ErratumD}.
For a local base parameter \(t\) vanishing at \(v\), smoothness and
invertibility of \(L\) give an exact sequence
\(0\to L\xrightarrow{t}L\to i_*L_v\to0\), where
\(i:X_v\hookrightarrow\mathcal X'\) is the closed embedding.
Its direct-image sequence gives, at \(v\),
\[
 0\longrightarrow F_v/tF_v\longrightarrow H^0(X_v,L_v)
 \longrightarrow (R^1\pi_*L)_v[t]\longrightarrow0.
\]
The last term denotes elements annihilated by \(t\). The support of
the torsion subsheaf of the coherent module \(R^1\pi_*L\) is closed
discrete, by the structure theorem for finitely generated modules
over the local discrete valuation rings of the smooth base curve.
Outside that support the canonical base-change map is an isomorphism.
Also \(F\) is torsion-free: multiplication by a nonzero base function
is injective on sections of \(L\). It is therefore locally free.
Its rank is one by choosing an interior point outside the discrete
torsion support. These statements concern the canonical map, not
merely equal dimensions.
Locally a frame gives a relative section \(\sigma\) of \(L\),
nowhere identically zero on a fibre. Its divisor is independent of the
frame. Let \(D\) be its reduced support. On \(V_0\), the divisor is
\(2S_v\) and \(D\) is exactly \(S_v\).

\par\noindent\textit{Step 3.} The map \(D\to V'\) is proper, since \(D\) is closed in a proper
space. There is no vertical divisor component: a vertical divisor in
the smooth total space would contain a whole irreducible fibre,
contrary to the nonzero restriction of \(\sigma\). Every horizontal
component has proper image all of the connected base. Since the
interior fibre is irreducible, there is exactly one horizontal
component, and its multiplicity in \(\operatorname{div}\sigma\) is
two. Thus \(\operatorname{div}\sigma=2D\) on this base. In the smooth
ambient space \(D\) is Cartier. It is flat over the curve because
its local rings are torsion-free over the base discrete valuation
ring. Restriction of the Cartier identity to an interior fibre gives
\(2D_v=2S_v\), whence \(D_v=S_v\) with its reduced structure.
Flatness and the smoothness of these fibres make \(D\to V'\) smooth
over \(V_0\). Its analytic nonsmooth locus therefore misses \(V_0\);
the proper mapping theorem once more makes its image discrete.
After deleting that image, \(D\to V^{\rm good}\) is a smooth proper
family of connected surfaces. Each deletion is closed discrete in the
then-current open curve. Their finite union is at most countable,
and the final good locus is open; we do not assert that successive
exceptional sets cannot accumulate at an earlier removed point.

\par\noindent\textit{Step 4.} Ehresmann's theorem \cite[\S1, Proposition, pp.~154--155]{Ehresmann52} identifies the smooth fibres along every path;
in particular their first Betti number is two. \cite[Theorem~11]{Buchdahl1999} then gives a K\"ahler metric on every fibre. Hence
\(h^{1,0}=1\), every holomorphic one-form is closed, and the integral
period map has a rank-two lattice image. Constant \(h^{1,0}\),
properness, and analytic base change supply a local relative
holomorphic one-form. Locally on a simply connected
base disc, a relative one-form and a basis of the free part of
\(H_1\) give two holomorphic periods. Their ratio lies in one
half-plane after orienting the basis. To see holomorphicity directly,
choose local primitives of the closed relative form on simply
connected fibre charts; their differences are holomorphic functions
of the base, and the finite sum around a transported cycle is its
period. The modular invariant of the ratio is independent of all
choices and is holomorphic. This proves the last assertion without
requiring a global choice of section of the Albanese family.
\end{proof}

\begin{theorem}[PUSH D.1.3; T]\label{push:D.1.3}
For every \(w_0\) with \(|w_0|=216\) and every connected open
neighbourhood \(U\) of \(w_0\), there is no flat proper holomorphic
map \(\mathcal X\to U\) from a reduced complex space whose fibres
over \(U\cap\{0<|w|<216\}\) are biholomorphic, fibre by fibre, to
\(X_w\). In particular an extension isomorphic to the given family
over that intersection does not exist.
\end{theorem}
\noindent\emph{Uses:} FAM A.1.1; PUSH D.1.1 and D.1.2.\par

\begin{proof}
We may restrict a putative extension to a small disc meeting the
circle, with closure away from zero. This is a holomorphic base
restriction; properness and flatness are preserved. The calculation
of the canonical evaluation in the earlier return, \ref{push:A.1.1}, gives
\(K_{X_w}=\mathcal O(-2S_2(w))\) and
\(h^0(K_{X_w}^{-1})=1\). Its sole divisor is \(2S_2(w)\), so it is
intrinsic under every fibrewise biholomorphism. The surface has
\(b_1=2\): on its finite product cover the rotation kills the moving
factor in invariant first cohomology and fixes the elliptic factor.
These are the exact FAM A.1.1 and FAM A.4.1 steps required in Lemma~\ref{push:D.1.2}.

The at most countable exceptional set in that lemma cannot contain
an arc of the circle. Choose a boundary point outside it and then a smaller disc
on which the reduced divisor is a smooth proper surface family. Its
holomorphic Albanese invariant agrees, pointwise on the inner part,
with \(J(iw/216)\). It is therefore an analytic continuation through
that point, contrary to Proposition~\ref{push:D.1.1}. No choice of the
fibrewise biholomorphisms was used. Thus the weaker hypothesis already
leads to the contradiction.
\end{proof}

\subsection{The specified pinching and its virtual Euler number}
\begin{definition}[PUSH A.1.setup]\label{push:A.1.setup}
We write $M_S=\mathcal O_S(P-O)$, $T=t_{6P}$,
$J_{\rm lin}=M_S\otimes T^*M_S^{-1}=\pi^*\mathcal O_B(1)$, and
$R=\mathbb P_S(\mathcal O\oplus M_S)$. The convention for this projective
bundle is that its zero and infinity sections have normal bundles
$M_S$ and $M_S^{-1}$. The identification defining $Q$ is
$s_0(x)\sim s_\infty(Tx)$. The equality defining $J_{\rm lin}$ is the equality
in \cite[Lemma~2.1]{E26}, with this direction of $T$; its intersection
calculation uses $P\cdot O=0$, $P^2=O^2=-1$,
$6P\cdot O=6P\cdot P=2$. Thus the degree on $O$ is
$(O-P)\cdot6P+(P-O)\cdot O=1$. We use the section
$\phi\in H^0(S,J_{\rm lin})$ whose divisor is $S_\infty$, as in \cite[Corollary~2.2]{E26}. 
\end{definition}
\noindent\emph{Uses:} E26 Lemma 2.1 and Corollary 2.2.\par

\begin{refuted}[E26 Conjecture 8.2; false as worded]\label{e26:Conjecture8.2}
There is a flat, proper filling $\mathcal X\to \Delta$ and the central fiber $X(0)$ is non-normal. The normalization
is $\mathbb P_S(\mathcal O\oplus M)$ and the normalization map glues the zero and infinity sections via $t_{6P}$.
\end{refuted}
\noindent\emph{Uses:} the untwisted family and the specified pinching of \cite[Conjecture~8.2]{E26}. In this quotation $M$ is Engel's bundle $M_S$.\par
\begin{lemma}[PUSH A.3.1; T]\label{push:A.3.1}
Let $Z$ be a compact reduced local-complete-intersection complex
space of pure dimension $n$. In any proper flat holomorphic
family over a disc with central fibre $Z$ and smooth nearby
fibres, their Euler number is
\[
 e_{\mathrm{vir}}(Z)=\langle c_n(T_Z^{\mathrm{vir}}),[Z]\rangle.
\]
This statement needs no projectivity or K\"ahler hypothesis.
\end{lemma}
\noindent\emph{Uses:} The local-complete-intersection deformation calculation; Levy Lemmas 3.1--3.2; \L{}ojasiewicz \S3 Theorems 1--2.\par

\begin{proof}
\par\noindent\textit{Step 1.} Flat deformation preserves the local-complete-intersection
property relative to the disc near $Z$. Locally this follows by
lifting an embedding and its regular sequence of equations;
Nakayama's lemma and flatness show that their lifts generate
the ideal and remain a regular sequence. Properness allows one
to shrink the disc so that these neighbourhoods cover the total
space. The relative cotangent complex is consequently perfect
of amplitude $[-1,0]$. Its dual defines a topological complex
$K$-class $V$ on that total space; local finite vector-bundle
complexes define the class by descent, or by stabilization and
partitions of unity; the analytic-to-topological comparison is
given in \cite[Lemmas 3.1--3.2 and Section 3.1]{Levy08}.
Its restriction to a smooth fibre is its
tangent bundle, and its restriction to $Z$ is its intrinsic
virtual tangent class.

\par\noindent\textit{Step 2.} The compact fundamental cycles of the central and a nearby
fibre are homologous in the total space. More explicitly,
the inverse image $K$ of a short closed radial segment
$[0,\epsilon]$ is compact and semianalytic. Away from
$\operatorname{Sing}Z$ it is an oriented real
$(2n+1)$-manifold with boundary the nearby fibre and
$-Z_{\mathrm{reg}}$, because flatness with smooth fibre
implies smoothness of the morphism at those points.
The exceptional set $\operatorname{Sing}Z$ has real
dimension at most $2n-2$. Triangulate compatibly with
this set and the two endpoints, using
\cite[Section 3, Theorems 1--2]{Loj64}. The sum of the
oriented top simplices has boundary exactly
$[X_\epsilon]-[Z]$: no real $2n$-simplex can be supported
on that exceptional set. Reducedness and flatness give
multiplicity one on every central component. Pairing the globally defined
$c_n(V)$ with these homologous cycles gives the equality of
Chern numbers. Gauss--Bonnet on the smooth compact complex
fibre gives its Euler number. This proves the analytic assertion
directly. The algebraic counterpart is \cite[Theorem~5.8 and Remark~5.9]{FG10}; those scheme statements alone would
not justify omitting the analytic comparison above.
\end{proof}

\begin{theorem}[PUSH A.3.2; T]\label{push:A.3.2}
For the specified pinching $Q$ one has $e_{\mathrm{vir}}(Q)=0$.
Every smooth proper-flat smoothing of $Q$ therefore has Euler
number zero. In particular there is no proper flat holomorphic
filling of Engel's untwisted family with central fibre exactly
$Q$, even if no map to $B$ is required. Thus Conjecture 8.2
of \cite[Conjecture~8.2]{E26} is false as worded.
\end{theorem}
\noindent\emph{Uses:} PUSH A.3.1; the specified pinching; the Euler number two of the chosen untwisted family.\par

\begin{proof}
\par\noindent\textit{Step 1.} Let $D_0,D_\infty\subset R$ be the conductor sections and
$\nu:R\to Q$ the normalization. On either branch of the
hypersurface $xy=0$, the derivative comparison with its
normalization gives in $K^0(R)$, where pullback of the
perfect tangent class means derived pullback,
\[
 \nu^*T_Q^{\mathrm{vir}}
 =T_R-i_{0*}(J_{\rm lin}|_{D_0})-i_{\infty*}(J_{\rm lin}|_{D_\infty}).
\]
Indeed in a smooth local ambient space the normal line $A$ of
the branch and the normal line of the whole hypersurface
differ by $D$: the latter is $A(D)$. Their difference in
$K$-theory is $-[i_*(A(D)|_D)]$, and $A(D)|_D$ is exactly
the product of the two branch normals, namely $\mathcal T^1_Q$.
These local comparisons are intrinsic and agree on overlaps.
On the infinity branch $J_{\rm lin}$ means its $T^{-1}$-pullback
from the double locus, canonically the same line bundle
because it comes from $B$. Thus $J_{\rm lin}$ extends from both sections by the pullback from $B$,
so $[i_*J_{\rm lin}]=[J_{\rm lin}]-[J_{\rm lin}(-D)]$.

\par\noindent\textit{Step 2.} Put $j=c_1(J_{\rm lin})=\pi_R^*F$, $d_i=[D_i]$. Since the sections are
disjoint, $d_0d_\infty=0$, and since $F$ is a fibre, $j^2=0$.
Consequently
\[
 c(\nu^*T_Q^{\mathrm{vir}})
 =c(T_R)\frac{(1+j-d_0)(1+j-d_\infty)}{(1+j)^2}
 =c(T_R)\{1-d_0-d_\infty+j(d_0+d_\infty)\}.
\]
The degree-three part has integral
\[
 e(R)-\sum_i\int_{D_i}c_2(T_R|_{D_i})
       +\sum_i\int_{D_i}j\,c_1(T_R|_{D_i}).
\]
We have $e(R)=2e(S)=24$. On a section with normal class
$n_i$, $c_2(T_R|_{D_i})=c_2(S)+c_1(S)n_i$ and
$c_1(T_R|_{D_i})=c_1(S)+n_i$. The two normal classes on
the unglued normalization are $m,-m$, where $m=c_1(M_S)$.
Their sum is zero; $c_1(S)=F$ and $F^2=0$. The displayed
number is therefore $24-24+0=0$. Because
$\nu_*[R]=[Q]$, this is the intrinsic virtual number on $Q$.
The projective-bundle calculation is independently checked
in \texttt{ab\_push\_local\_checks.py}, using
$H(H+m)=0$, $d_0=H+m$, $d_\infty=H$.

\par\noindent\textit{Step 3.} The untwisted quotient away from its cusp has Euler number zero:
scalar multiplication supplies a free circle action, including
the two additive places before logarithmic transformation. Indeed
a circle scalar identified with an iterate of $u\phi$ must be
the identity, since every positive iterate strictly contracts
the Hermitian norm and every negative iterate expands it.
Compact fibres with a free circle action have Euler number zero.
The chosen Mumford fibre has two zero-dimensional torus orbits
and all other torus strata have Euler number zero, so
$e(X(u))=2$. The same number follows independently from the
semistable nearby-cycle count in PUSH A.4.3 in the repository. Lemma~\ref{push:A.3.1}
therefore excludes such a smoothing of $Q$. No requirement
on its leading smoothing section, and no extension of $f$,
entered this argument.

\par\noindent\textit{Step 4.} There is also a direct check that zero is the correct smooth
number. Replace $\phi$ in the global construction of \ref{push:A.1.2}
by a section of $J_{\rm lin}$ vanishing on a smooth elliptic fibre of
$S$. The compact collar construction still applies. At its
zero divisor the fixed nonzero-parameter equation is $xy=uz$
with $dz\ne0$, so every nearby threefold is smooth. The
double-locus nearby Euler function is one on that elliptic
curve and zero off it; its integral is zero. This is an actual
proper-flat smoothing of $Q$, but is not Engel's family.
\end{proof}

\begin{announcement}[PUSH A.3.2, final sentence; A]\label{push:A.3.2.replacement}
The proper nodal
filling of \ref{push:A.1.2} and the two natural semistable fillings of
\ref{push:A.2.2} are precise replacements, with different central fibres.
\end{announcement}
\noindent\emph{Uses:} PUSH A.1.2 and A.2.2.\par
The mechanism changes the central fibre through the nodal construction or one of the two toroidal resolutions described below.

\noindent\emph{Proof:} proof in the repository, PUSH A.1.2 and A.2.2 (\url{https://github.com/KokunoYumeto/s6-complex-structure-workbench}).
\paragraph{Counterexample checklist.}
The target is the exact pinching $Q$ in the quotation. Its normalization
and gluing are fixed in Definition~\ref{push:A.1.setup}; its reduced local
complete-intersection structure is the local equation $xy=0$ used in
the proof. Statement~\ref{push:A.3.1} applies to every proper flat smooth
smoothing of that $Q$, without a projectivity or K\"ahler assumption.
Statement~\ref{push:A.3.2} gives the conflicting counts $0$ and $2$.
The announced replacements change the central fibre, and the nodal
replacement also changes the nonzero fibres by the specified contraction.
\subsection{Announced corrected fillings}
\begin{definition}[PUSH A.1.fan]\label{push:A.1.fan}
At fixed $h_{t_c},h_u\geq0$, the cone $P_{k,l}$ cuts out the rectangle
\[
 kh_{t_c}\leq y_1\leq(k+1)h_{t_c},\qquad
 l(h_{t_c}+h_u)\leq y_2\leq(l+1)(h_{t_c}+h_u).
\]
\end{definition}
\noindent\emph{Uses:} The cusp cocharacter lattice.\par

\begin{announcement}[PUSH A.1.1; A]\label{push:A.1.1}
The cones $P_{k,l}$ in Definition~\ref{push:A.1.fan} and their faces form a fan. Its two
period automorphisms on the cocharacter lattice are
\[
 G_1(y_1,y_2,h_{t_c},h_u)=(y_1+h_{t_c},y_2,h_{t_c},h_u),\qquad
 G_2(y_1,y_2,h_{t_c},h_u)=(y_1,y_2+h_{t_c}+h_u,h_{t_c},h_u).
\]
The affine chart of $P_{0,0}$ is
\[
 U_{0,0}=\operatorname{Spec}\mathbb C[x,y,z,v,u]/(xy-uzv),
 \qquad t_c=zv.
\]
In particular the unsubdivided chart is not a smooth simultaneous
model of the cusp.
\end{announcement}
\noindent\emph{Uses:} PUSH A.1.fan.\par

The mechanism is the rectangle tiling in the cocharacter lattice and the dual-monoid equation $xy=ut_c$.

\noindent\emph{Proof:} proof in the repository, PUSH A.1.1 (\url{https://github.com/KokunoYumeto/s6-complex-structure-workbench}).

\begin{announcement}[PUSH A.1.2; A]\label{push:A.1.2}
There is a reduced complex space $\mathcal X^\flat$ and a proper flat
holomorphic map
\[
 (f,\varpi):\mathcal X^\flat\longrightarrow B\times\Delta_u
\]
whose central threefold is $Q$. For $u\ne0$ its fibre is the
contraction of Engel's chosen $X(u)$ at one rational curve of the
Mumford fibre. This curve is the double curve corresponding to the
diagonal of the square in the chosen $A_2$ triangulation; in the
return's coordinates its direction is $e_1-e_2$. Its normal bundle
in $X(u)$ is $\mathcal O_{\mathbb P^1}(-1)^{\oplus2}$. Its image is
the unique ordinary double point of $\mathcal X^\flat_u$.
\end{announcement}
\noindent\emph{Uses:} PUSH A.1.1; E26 Lemma 2.1 and Corollary 2.2.\par

The mechanism glues a compact annular core to the global node bundle $a\otimes b=u\phi_x$, using the two specified collars. At the nodal elliptic fibre, the chart $xy=uzv$ produces the stated ordinary double point and its prescribed contraction.

\noindent\emph{Proof:} proof in the repository, PUSH A.1.2 (\url{https://github.com/KokunoYumeto/s6-complex-structure-workbench}).

\begin{definition}[PUSH A.2.rays]\label{push:A.2.rays}
Write the six rays of $P_{0,0}$ as
$A=(0,0,1,0)$, $B=(1,0,1,0)$, $C=(0,1,1,0)$,
$D=(1,1,1,0)$, $E=(0,0,0,1)$, $F=(0,1,0,1)$.
The return's $A_2$ diagonal is $BC$.
\end{definition}
\noindent\emph{Uses:} PUSH A.1.1.\par

\begin{announcement}[PUSH A.2.1; A]\label{push:A.2.1}
Among subdivisions with no new rays, the fan in \ref{push:A.1.1} has exactly
two smooth period-invariant subdivisions restricting to the chosen
$A_2$ triangulation at $h_u=0$. These are also all subdivisions
which are flat over the original $(t_c,u)$ bidisc, have reduced
$u$-central divisor, and retain that boundary triangulation.
Without the joint-flatness restriction there are infinitely many
smooth invariant subdivisions retaining that triangulation and
having reduced $u$-central divisor.
\end{announcement}
\noindent\emph{Uses:} PUSH A.1.1 and the retained bottom triangulation.\par

The mechanism enumerates the unimodular triangulations of the six-ray prism and imposes the period identifications. Joint flatness excludes a ray over the origin of the two-dimensional base; without that restriction, the stated period-invariant point stars give arbitrarily many refinements.

\noindent\emph{Proof:} proof in the repository, PUSH A.2.1 (\url{https://github.com/KokunoYumeto/s6-complex-structure-workbench}).

\begin{announcement}[PUSH A.2.2; A]\label{push:A.2.2}
The two minimal subdivisions of \ref{push:A.2.1} give proper flat holomorphic
families $\mathcal X^+,\mathcal X^-\to\Delta_u$ with smooth total
spaces. Their joint maps to $B\times\Delta_u$ are flat. Over
$\Delta_u^*$ they are isomorphic as families to the specified
untwisted family $X(u)$. Each central fibre $Q^\pm$ is reduced,
irreducible, and has normal crossings, with smooth double locus
$S$ and no triple points.

Its normalization $R^\pm$ is obtained by blowing up $R$ along
$S_\infty$ in, respectively, one of its two sections and then
taking the small resolution at the unique ordinary double point
of that blow-up prescribed by the fixed bottom diagonal.
The two proper transforms of the sections are still copies of
$S$, glued by $T$. Their normal-bundle product is trivial;
in particular $\mathcal T^1_{Q^\pm}\simeq\mathcal O_S$.
\end{announcement}
\noindent\emph{Uses:} PUSH A.1.2 and A.2.1.\par

The mechanism applies either of the two period-invariant subdivisions to the proper nodal construction and leaves the other charts unchanged. The central heights determine the reduced normal crossings, and the conductor normal-bundle calculation gives the trivial product.

\noindent\emph{Proof:} proof in the repository, PUSH A.2.2 (\url{https://github.com/KokunoYumeto/s6-complex-structure-workbench}).

\begin{announcement}[PUSH A.2.3; A]\label{push:A.2.3}
The unrestricted request that every normalized component
be a blow-up of $R$ is false. The first maximal-cone
star already contributes a $\mathbb P^3$ component,
which admits no proper bimeromorphic map to $R$, since
such a map would inject a rank-eleven rational $H^2$
into a rank-one group. \ref{push:A.2.4} gives every normalization
by explicit star charts; A.2.5 supplies the terminating
blow-up and blow-down algorithm with centres, order,
exceptional bundles and gluing maps. Thus the literal
shape is refuted and its general algorithmic
replacement is proved.
\end{announcement}
\noindent\emph{Uses:} PUSH A.2.4 and A.2.5.\par

The mechanism is the new complete toric component in a maximal-cone star. Its second Betti number obstructs the proposed bimeromorphic map to the ruled threefold; the finite star-chart description is retained in its place.

\noindent\emph{Proof:} proof in the repository, PUSH A.2.3 (\url{https://github.com/KokunoYumeto/s6-complex-structure-workbench}).

\begin{definition}[PUSH A.2.periods]\label{push:A.2.periods}
Write the period shears as
\begin{align*}
 G_1(a,b,k,\ell)&=(a+k,b,k,\ell),\\
 G_2(a,b,k,\ell)&=(a,b+k+n\ell,k,\ell).
\end{align*}

In a unimodular chart its
coordinates $z_i$ satisfy
\[
 r=\text{unit}\cdot\prod_{i=1}^4 z_i^{\ell_i},
 \qquad t_c=\text{unit}\cdot\prod_{i=1}^4 z_i^{k_i}.
 \tag{A.24d}
\]
\end{definition}
\noindent\emph{Uses:} PUSH A.2.4, finite fan coordinates.\par

\begin{announcement}[PUSH A.2.4; A]\label{push:A.2.4}
The following finite-data procedure parametrizes all smooth
period-invariant subdivisions of the cusp fan after
$u=r^n$, retaining the prescribed bottom triangulation and
having finitely many cones modulo the periods. Here $n=1$
gives PUSH Part A, and $n=2,12$ give the two cases of PUSH B.2.
Each output defines a proper flat family over $\Delta_r$
with smooth total space, whose punctured restriction is the
given untwisted family pulled back by $u=r^n$.

For a bound $H\geq1$, retain the four bottom rays
$A,B,C,D$ and consider all primitive vectors
\[
 \begin{gathered}
 v=(a,b,k,\ell),\qquad
 \ell\geq1,\quad k\geq0,\quad k+\ell\leq H,\\
 0\leq a\leq k,\qquad0\leq b\leq k+n\ell.
 \end{gathered}
 \tag{A.24a}
\]
Take any collection of four-element subsets of this finite
set satisfying the following exact tests:
\begin{enumerate}
\item Each four-column determinant has absolute value one.
\item After dividing each ray by $k+\ell$, the resulting
tetrahedra intersect pairwise in common faces, and their
volumes sum to $(n+2)/6$.
\item Their bottom boundary is the two triangles $ABC,BCD$.
Their boundary subdivisions on $a=0$ and $a=k$ agree under
$G_1$, and those on $b=0$ and $b=k+n\ell$ agree under $G_2$.
\end{enumerate}
All these are finite rational-arithmetic tests. Extend a
passing collection by the two periods, and discard duplicates
as $H$ increases. This is exhaustive. For a reduced central
divisor, retain only $\ell=1$ in (A.24a).

The central prime components are indexed by period orbits of
the selected rays with $\ell>0$, and the multiplicity of the
component indexed by $v$ is exactly $\ell$. The normalization
charts of that component are the star fan of $v$ in
$N/\mathbb Zv$, with the inherited bundle units and the
stabilizer action. For $k>0$ the stabilizer is trivial and this
normalization is a smooth complete toric threefold. All
multiple loci, including possible triple and quadruple loci,
are specified by the positive-$\ell$ rays belonging to each
cone. In the reduced case
$\mathcal T^1\simeq\mathcal O_{\operatorname{Sing}D}$,
so every output is d-semistable.
\end{announcement}
\noindent\emph{Uses:} PUSH A.1.1 and the explicit base-changed period shears.\par

The mechanism enumerates a finite height-bounded collection of primitive rays, tests unimodularity and rational face intersections, and imposes the two period identifications. Increasing the bound exhausts the finite invariant subdivisions, while the star fans record the central normalizations and their multiplicities.

\noindent\emph{Proof:} proof in the repository, PUSH A.2.4 (\url{https://github.com/KokunoYumeto/s6-complex-structure-workbench}).

\begin{announcement}[PUSH A.2.5; A]\label{push:A.2.5}
The requirement that every normalized component of every
smooth refinement be a blow-up of
$R=\mathbb P_S(\mathcal O\oplus M_S)$ is false, already
for a reduced refinement with $n=1$.
The following is a uniform replacement for any finite
fan certificate $\Sigma$ in \ref{push:A.2.4}.

The star charts and integer transition matrices of \ref{push:A.2.4}
give every normalized component, its gluing maps, and
all multiple loci directly. There is also a terminating
finite-data algorithm that expresses its comparison
with a reference model as an explicit sequence of
blow-ups and blow-downs along specified smooth strata.
At each blow-up the algorithm gives the normalized
strict transforms, the exceptional projective bundle,
its multiplicity, and every new intersection.
This algorithm allows blow-downs; no toric
strong-factorization assertion is used.

For a component of multiplicity $\ell$, the local
nonreduced central equation is the monomial in A.24d.
More generally, in a unimodular chart with central
equation $F=\prod_i z_i^{\ell_i}$, the cotangent-complex
first-order deformation sheaf is
\[
 \mathcal T^1_D
  =\mathcal O_D\Big/
    \left(\ell_i z_i^{\ell_i-1}
          \prod_{j\ne i}z_j^{\ell_j}:\ell_i>0\right).
 \tag{A.25a}
\]
Here $\mathcal T^1_D$ means
$\mathcal Ext^1(\mathbb L_D,\mathcal O_D)$; in the
reduced case it also equals
$\mathcal Ext^1(\Omega_D,\mathcal O_D)$.
The reduced case is the trivial line on the singular
scheme, as in \ref{push:A.2.4}. For a nonreduced centre the displayed
Jacobian quotient is the appropriate deformation sheaf;
the reduced normal-crossings notion of d-semistability
is not asserted for it.
\end{announcement}
\noindent\emph{Uses:} PUSH A.2.4 and its marked star charts.\par

The mechanism compares finite marked cone complexes through ordinary star blow-ups and their inverse moves. The local monomial Jacobian quotient records the deformation sheaf with its stated nonreduced scope.

\noindent\emph{Proof:} proof in the repository, PUSH A.2.5 (\url{https://github.com/KokunoYumeto/s6-complex-structure-workbench}).

\section{Announcements at the puncture and in the space of structures}
\label{sec:announcements}
\subsection{The fixed-family noncompletion statement}
\begin{definition}[PUSH B.4.33.scope]\label{push:B.4.33.scope}
The retained sphere family has multiplicities $(3,4)$ and the fixed data
$\epsilon=(1,2,-4,0)^t$, $\epsilon'=(1,3,-3,0)^t$,
$v_1=\epsilon$, $v_2=-\epsilon'$, $l_0=0$, and $p=12l_0-4l_1-3l_2=-1$.
The untwisted family is
$X(u)=M_S^{\times,\circ}/\langle u\phi\rangle$, completed at infinity
by its Mumford fibre; the sphere family is $X_w$.
\end{definition}
\noindent\emph{Uses:} The first-run retained data.\par

\begin{announcement}[PUSH B.4.33; A]\label{push:B.4.33}
The actual punctured six-sphere family has no proper
Hausdorff complex-analytic filling after any finite
cover of its degeneration parameter. In particular,
there is no positive integer $N$ for which the cover
$w=r^N$ has a proper filling with smooth total space
and reduced normal-crossings central fibre.

The assertion permits an arbitrary singular central
space, arbitrary additional components, and arbitrary
central multiplicities. It does not impose a
K\"ahler, projective, or logarithmic hypothesis on
the putative filling. Its proof derives the specific
logarithmic torsor needed for the contradiction from
differences of sections extending into that filling.
\end{announcement}
\noindent\emph{Uses:} PUSH B.4.31 and B.4.32; the actual fixed punctured sphere family.\par

The mechanism passes from a proposed filling to a marked proper map and controls the radial-loop class through differences of extending sections. The resulting specific logarithmic torsor contradicts the nonzero degree; an arbitrary affine origin is not asserted to extend.

\noindent\emph{Proof:} proof in the repository, PUSH B.4.33 (\url{https://github.com/KokunoYumeto/s6-complex-structure-workbench}).

The assertion has the fixed-family scope just stated. Whether the same
noncompletion holds for every admissible twist remains a question.
The proposed reading of a monodromy integer as $J=P\mathbin{.}O$ is a
reader's remark recorded in the intake; it is unverified and is not a
theorem or a premise of this paper.
\subsection{Smooth monodromy: two runs and two methods}
\begin{definition}[PUSH C.orientation]\label{push:C.orientation}
The six-sphere has its complex orientation and a positive $w$-circle
is counterclockwise. We use
$M_\Phi=(S^6\times[0,1])/((x,1)\sim(\Phi(x),0))$ with orientation
$dt\wedge\operatorname{or}(S^6)$, and
$\Sigma_\Phi=D^7\cup_\Phi(-D^7)$ for the identification with
$\Theta_7$; its positive generator has Eells--Kuiper invariant $1/28$.
\end{definition}
\noindent\emph{Uses:} The complex orientation and the positive parameter circle.\par

\begin{definition}[PUSH C.4.1.scope]\label{push:C.4.1.scope}
The twist comparison used in PUSH C.3.24 has fixed multiplicities $(3,4)$ and the admissible sphere data $|p|=1$ of FAM B.2.1. Its monodromy convention is the positive $w$-circle with the complex orientation on every fibre.
\end{definition}
\noindent\emph{Uses:} FAM B.2.1; PUSH C.4.1.\par

\begin{announcement}[PUSH C.3.24; A]\label{push:C.3.24}
For the actual retained $S^6$ family and a positive generator
of its punctured $w$-parameter disc, the oriented smooth
monodromy satisfies
\[
 [\Phi]=0\quad\hbox{in }\pi_0\operatorname{Diff}^+(S^6)
                       \simeq\Theta_7\simeq\Z/28.
\]
Its mapping torus is consequently oriented-diffeomorphic to
$S^1\times S^6$. Every finite cyclic base change has trivial
smooth monodromy as well. By C.4.1's marked family comparison,
the same conclusion applies to the admissible sphere twist
data covered by that comparison.
\end{announcement}
\noindent\emph{Uses:} PUSH C.3.18, C.3.21--C.3.23 and C.4.1.\par

The mechanism is an integer calculation before reduction modulo $28$: the source gives $p_{\rm rel}=-24$, $s_\delta=-2$ and $J_\psi=\epsilon_{\rm cap}(-24-12(-2))/24=0$. Additivity compares the six-loop representative to six times the one-loop integer.

\noindent\emph{Proof:} proof in the repository, PUSH C.3.24 (\url{https://github.com/KokunoYumeto/s6-complex-structure-workbench}).

\begin{announcement}[PUSH C.3.27; A]\label{push:C.3.27}
The actual positive-loop smooth monodromy satisfies
$[\Phi]=0\in\Theta_7$. This determination uses the explicit
eight-dimensional coboundary of C.3.25 and the direct index
additivity of C.3.26. It does not use the three-knot integer,
Wall's formula for it, the Davis generator comparison, or the
conclusion of \ref{push:C.3.24}.
\end{announcement}
\noindent\emph{Uses:} PUSH C.3.25 and C.3.26; the shared characteristic calculations C.3.18--C.3.22.\par

The mechanism uses the explicit eight-dimensional spin$^c$ coboundary and the integer index $\lambda=(p_\psi-4k_\psi)/24$, with $p_\psi=-24\epsilon_{\rm cap}$ and $k_\psi=-6\epsilon_{\rm cap}$. The boundary-invariant and one-loop arguments differ from the orbit-knot method; the actual annular representative and characteristic inputs are shared.

\noindent\emph{Proof:} proof in the repository, PUSH C.3.27 (\url{https://github.com/KokunoYumeto/s6-complex-structure-workbench}).

\begin{announcement}[MAPS B.6.21; A]\label{maps:B.6.21}
For the original $(3,4)$ family, with its complex orientation,
the positive $w$ loop and the generator of $\Theta_7$ fixed by
$\mu=1/28$, the monodromy satisfies
\[
                         [\Phi]=0\in\mathbb Z/28.
 \tag{B.6.38}
\]
It is therefore smoothly isotopic to the identity.  Its sixfold
class is zero, and among the two resulting K7 residues $0,14$
the original class is $0$.  This statement does not yet assert
the same value for the other multiplicity pairs.
\end{announcement}
\noindent\emph{Uses:} MAPS B.6.15--B.6.20 and the specified mapping-torus comparison.\par

The mechanism uses the free-circle replacement with its actual peripheral identifications, the values $k=-1$ and $\ell=-4$, and the parametrized spin-surgery trace. Its normalized correction is zero, so the specified Eells--Kuiper invariant selects the original residue $0$.

\noindent\emph{Proof:} proof in the repository, MAPS B.6.21 (\url{https://github.com/KokunoYumeto/s6-complex-structure-workbench}).

\begin{definition}[PUSH E.1.1.scope]\label{push:E.1.1.scope}
For the higher-family announcements, $m=3k$, $n=4l$,
$\gcd(m,n)=1$, and $k,l>0$. A fixed sphere twist datum in the K1 scope
of MAPS means fixed integers $\rho_0,N_0,\rho_1,N_1$ with
\[
 w_1=N_0\varepsilon-\rho_0\widehat\delta,\qquad
 w_2=N_1\varepsilon'-(\rho_1+N_1)\widehat\delta,\qquad l_0=0,
\]
and $nN_0+mN_1=\pm1$, in the admissible period domain. The vectors are
the first-run clockwise-root vectors; the original K1 data use the
transported counterclockwise-root marking. The parameter period is
$2kl\mathbb Z$.
\end{definition}
\noindent\emph{Uses:} PUSH E.1.1--E.1.3; FAM B.3.1.\par

\begin{announcement}[MAPS B.6.24; A]\label{maps:B.6.24}
For every admissible coprime sphere pair $(3k,4l)$ and every
fixed twist datum of K1, the monodromy for the positive
generator $2kl$ of its parameter-period lattice satisfies
\[
                    [\Phi_{3k,4l}]=0\in\Theta_7.
 \tag{B.6.43}
\]
The identifications and orientation are those specified in
the question.  In particular this completes the numerical
part of B.6 for both the original and higher families.
\end{announcement}
\noindent\emph{Uses:} The fixed K1 twist datum; MAPS B.6.19 and B.6.22--B.6.23.\par

The mechanism is the higher-family free-circle quotient calculation $k_{\rm cup}=-kl$, $\ell=-4kl$, followed by the signed spin-surgery trace and the mapping-torus comparison. This is the source statement for every fixed K1 datum in the stated scope.

\noindent\emph{Proof:} proof in the repository, MAPS B.6.24 (\url{https://github.com/KokunoYumeto/s6-complex-structure-workbench}).

The last announcement has a broader scope than MAPS B.6.21, whose
last sentence is retained. The supplied open-problem list also calls
``the monodromy of the other families'' open. Those wordings are not
silently identified: MAPS B.6.24 announces zero for the explicitly
specified coprime sphere pairs and fixed K1 data, while an open question
beyond that scope is not answered here.
\subsection{Families and components of the space of structures}
\begin{definition}[PUSH E.2.5.scope]\label{push:E.2.5.scope}
The smooth-completion hypotheses used below are those of PUSH E.2.5:
$\pi:\mathcal X\to\Delta$ is a proper holomorphic submersion,
all fibres are diffeomorphic to $S^6$, and every noncentral fibre belongs
to one fixed coprime multiplicity family $(3k,4l)$ of PUSH E.2.1.
The integers in that statement are $A=12kl$ and $D=kl-3k-4l\ne0$.
For fixed discrete logarithmic data, the isomorphism parameter of
PUSH E.2.7 is denoted here by
\[
 w_{k,l}=\exp\!\left(-\frac{\pi i c}{kl}\right).
\]
The symbol $z$ in those source statements denotes this parameter,
not the uniformizing coordinate of the first-run conventions.
\end{definition}
\noindent\emph{Uses:} PUSH E.2.5 and E.2.7.\par

\begin{announcement}[PUSH E.2.21; A]\label{push:E.2.21}
Let $\pi:\mathcal X\to\Delta$ satisfy the hypotheses of
Definition~\ref{push:E.2.5.scope}, with fixed coprime multiplicities $(3k,4l)$ on
every noncentral fibre.  Then $X_0$ is a member of that
same constructed family.  Its isomorphism parameter is
the interior value $w_{k,l}(0)$ of PUSH E.2.7.

After choosing a holomorphic lift $c(s)$ of that parameter
and shrinking the disc, the actual family is isomorphic
to the prescribed family at $c(s)$.  The ambiguity in
this comparison is the connected translation automorphism
group already determined in E.1.
\end{announcement}
\noindent\emph{Uses:} PUSH E.2.5, E.2.7 and E.2.10--E.2.20.\par

The mechanism extends the intrinsic pencil and its nonzero interior period, identifies the ordinary marked Jacobian, and glues the finite-root and cusp comparisons using the connected translation group.

\noindent\emph{Proof:} proof in the repository, PUSH E.2.21 (\url{https://github.com/KokunoYumeto/s6-complex-structure-workbench}).

\begin{announcement}[PUSH E.2.22; A]\label{push:E.2.22}
Let $\pi:\mathcal X\to\Delta$ be a proper holomorphic
submersion with fibres diffeomorphic to $S^6$.  Suppose
there is a nonempty open subset $U\subset\Delta$ such
that every $X_s$, $s\in U$, is a member of the fixed
coprime multiplicity family $(3k,4l)$ of E.2.1.
Then every fibre over $\Delta$ is a member of that
same family.
\end{announcement}
\noindent\emph{Uses:} PUSH E.2.1 and E.2.21.\par

The mechanism propagates the geometric data outside successive discrete analytic exceptional sets and then fills those sets in reverse order by the smooth-completion comparison.

\noindent\emph{Proof:} proof in the repository, PUSH E.2.22 (\url{https://github.com/KokunoYumeto/s6-complex-structure-workbench}).

\begin{announcement}[PUSH E.2.23; A]\label{push:E.2.23}
In a proper smooth holomorphic family of compact complex
manifolds over a connected reduced complex analytic base,
if one fibre is a sphere member of E.2.1 with coprime
multiplicities $(3k,4l)$, then every fibre is a constructed
member with that same pair.

Consequently two of the constructed sphere members are
deformation equivalent, allowing chains of proper smooth
families over connected complex analytic bases, if and
only if their ordered multiplicity pairs agree.  In
particular the distinct coprime pairs give distinct
deformation classes.
\end{announcement}
\noindent\emph{Uses:} PUSH E.2.21 and E.2.22; local miniversality.\par

The mechanism combines propagation on Riemann surfaces with the local analytic-base argument. It uses the constructed fixed-pair locus and does not assume that a punctured analytic disc can be selected inside an arbitrary open set.

\noindent\emph{Proof:} proof in the repository, PUSH E.2.23 (\url{https://github.com/KokunoYumeto/s6-complex-structure-workbench}).

\begin{announcement}[PUSH E.2.24; A]\label{push:E.2.24}
Let $\mathcal I(S^6)$ be the space of smooth integrable
almost-complex structures on the fixed smooth sphere,
with its $C^\infty$ topology.  For a fixed coprime pair
$(3k,4l)$, let $\mathcal I_{k,l}$ consist of the structures
whose complex manifolds are biholomorphic to constructed
members with that pair.  Then $\mathcal I_{k,l}$ is open
and closed in $\mathcal I(S^6)$.

Consequently a continuous family of such structures
over any connected topological parameter space cannot
contain constructed members with two different pairs.
The same separation holds in the quotient spaces by
$\operatorname{Diff}_0(S^6)$ or by all orientation-preserving
diffeomorphisms, with their quotient topologies.
This statement does not identify the number of marked
connected components belonging to a single pair.
\end{announcement}
\noindent\emph{Uses:} PUSH E.2.23 and local completeness of the deformation charts.\par

The mechanism transfers the analytic-family separation to the local deformation charts in the smooth topology of \cite[\S2.1 and \S3.2, Theorem~3.4 and Remark~3.6]{Meersseman2019E}, and then to the indicated quotient spaces. It supplies no count of marked connected components within a single pair.

\noindent\emph{Proof:} proof in the repository, PUSH E.2.24 (\url{https://github.com/KokunoYumeto/s6-complex-structure-workbench}).

